%% ═══════════════════════════════════════════════════════════════════════════
%% CAPÍTULO 4 — METODOLOGIA   (versão da autora, 12/11/2026)
%%
%% CONVENÇÃO DE COR (definida em 000-main.tex):
%%   \revisao{}  = PRETO — rodada anterior, já incorporada
%%   \corr{}     = AZUL  — correções desta rodada
%%   \pendente{} = LARANJA — pendência que exige ação antes da entrega
%%
%% ─── PENDÊNCIAS QUE RESTAM (atualizado 12/11/2026) ─────────────────────────
%%  #8  [RESOLVIDA 12/11/2026] secao/pagina da faixa de temperatura: Secao 1.1, Tabela 1, p. 7; conversao: Secao 5.3.7, p. 28
%%       (conferir tambem: citacao DS000-19 vs DS001 atual da Bosch)
%%  #9  [RESOLVIDA 12/11/2026] RoboClaw #29785 e batmon #29787 reportados; sumd e MicroStrain nao reportados; ESBMC #8524
%%       RoboClaw, batmon) .....................
%% FECHADAS 12/11/2026: #1 (opções do workflow -> esbmc-verify.yml, lido no fork),
%%  #2 (ref da declaração de IA), #3 (replay FIFO/GPS -> limitação),
%%  #4 (harness GPS reexecutado: ESBMC 8.2.0/Z3, ~389s, SUCCESSFUL), #5 (ArduPilot
%%  opcional), #6 (CI = smoke check CBMC; tempo descartado), #7 (lista SAGRES ->
%%  limitação), #8-linhas (expo l.84, send_bytes l.852, FIFO SIZE l.152,
%%  gps.cpp l.741 conferido).
%% ───────────────────────────────────────────────────────────────────────────
%% ═══════════════════════════════════════════════════════════════════════════

\chapter{Metodologia}
\label{cap:metodologia}

%% LC 01 (resolvido 19/10/2026): "Incluir tabela: ferramentas, versoes, flags e
%% hipoteses de modelagem para cada etapa." -> Tabela tab:etapas-ferramentas.
%% [RESOLVIDO] Metodologia promovida a oito etapas. O replay e a validacao
%% cruzada, antes implicitos na Etapa 5, passam a ser as Etapas 6 e 7.

Este capítulo apresenta a metodologia empregada na verificação formal de módulos críticos do PX4 Autopilot com o ESBMC~\cite{Menezes2024}. O procedimento está organizado em oito etapas, conforme a numeração adotada no artigo derivado: (i) seleção dos módulos críticos; (ii) instrumentação para abstrair dependências de \textit{hardware}; (iii) especificação das propriedades de segurança; (iv) configuração e execução do ESBMC; (v) análise dos contraexemplos e conferência com o código de produção; (vi) reporte à comunidade PX4; (vii) validação cruzada com o CBMC e verificação de \textit{patches}; e (viii) reprodução em outro autopiloto (ArduPilot) e integração contínua. O módulo GPS uBlox é utilizado como exemplo para apresentar, de forma concreta, a aplicação dessas etapas.

\corr{A metodologia foi aplicada em duas campanhas, que totalizam 22 módulos distintos. A \textbf{primeira campanha}, apresentada no artigo derivado, abrangeu doze módulos relacionados a sensoriamento, estimação, navegação e comunicação. Nessa etapa, os \textit{harnesses} foram construídos a partir da transcrição de trechos do código do PX4. A \textbf{segunda campanha}, desenvolvida no âmbito desta dissertação, acrescentou dez módulos em três frentes: oito candidatos adicionais em módulos de entrada e de sistema, o driver de saída RoboClaw e o driver de bateria batmon. Nessa segunda etapa, a verificação e a conferência foram realizadas sobre o código-fonte \emph{original} do PX4, com \textit{replay} sob sanitizadores (Seção~\ref{sec:segunda-campanha}). As duas campanhas utilizaram o commit \texttt{8be2590c15}, exceto o batmon, cuja análise foi conferida contra o \texttt{main} de 13/10/2026 (commit \texttt{c9b9e4e25e}).}

As etapas~(i) e~(iv)--(vi) seguem o fluxo de Sousa et
al.~\cite{Sousa2026}; as etapas~(ii) e~(iii) são específicas de \textit{firmware},
e as etapas~(vii) e~(viii) são acréscimos deste trabalho. Dentro da etapa~(v), a
conferência contra o código de produção é a que separa violação alcançável no
modelo de falso positivo: foi ela que reclassificou o candidato do contador de
FIFO do BMI088 (Seção~\ref{subsec:imu-fifo}), que já havia sido proposto
\textit{upstream}, como falso positivo.

%% LC 02 (resolvido 19/10/2026): ordem real da primeira campanha.
\corr{Embora as etapas sejam apresentadas em ordem lógica, sua execução não seguiu essa sequência em todos os casos da primeira campanha. Nos casos do GPS e do contador de FIFO do BMI088, o reporte (vi) ocorreu imediatamente após a execução do ESBMC (iv), pois as \textit{issues} foram abertas para solicitar a triagem da comunidade. A conferência com o código de produção (v) foi realizada posteriormente, a partir da análise do código e das manifestações dos mantenedores. Assim, nesses casos, a sequência efetiva foi (i)--(iv), (vi), (v) e (vii). Essa inversão é considerada na discussão das ameaças à validade (Seção~\ref{sec:ameacas-validade}). Na segunda campanha, adotou-se a ordem prevista: nenhum candidato foi reportado antes da conferência com o código de produção e da execução do \textit{replay} (Seção~\ref{sec:segunda-campanha}).}

\corr{A Tabela~\ref{tab:etapas-ferramentas} resume, para cada etapa, as
ferramentas, versões, opções e hipóteses de modelagem empregadas.}

\begin{table}[!htbp]
\centering
\caption{\corr{Ferramentas, versões, opções e hipóteses de modelagem por etapa}}
\label{tab:etapas-ferramentas}
\scriptsize
\renewcommand{\arraystretch}{1.1}
\setlength{\tabcolsep}{3pt}
\newcommand{\cel}{\raggedright\hspace{0pt}}
\begin{tabular}{|p{2.0cm}|p{3.7cm}|p{4.9cm}|p{4.1cm}|}
\hline
\cel\textbf{Etapa} & \cel\textbf{Ferramentas e versões} & \cel\textbf{Opções} & \cel\textbf{Hipóteses de modelagem} \tabularnewline
\hline
\cel\corr{(i) Seleção} & \cel\corr{leitura do código e da documentação do PX4 (commit \texttt{8be2590c15})} & \cel\corr{---} & \cel\corr{módulo isolável por \textit{stubs}; amostra não estatística} \tabularnewline
\hline
\cel\corr{(ii) \textit{Stubs}} & \cel\corr{C/C++; \texttt{nondet\_T()}, \texttt{\_\_ESBMC\_assume} (1.ª campanha) e \texttt{\_\_VERIFIER\_assume} (2.ª)} & \cel\corr{---} & \cel\corr{E/S de \textit{hardware} sobreaproximada por valores não determinísticos; pré-condições da Tabela~\ref{tab:hipoteses-harness}} \tabularnewline
\hline
\cel\corr{(iii) Propriedades} & \cel\corr{\textit{assertions} no \textit{harness}} & \cel\corr{---} & \cel\corr{tradução de documento (\textit{datasheet}, manual, código) em predicado sobre variável de programa} \tabularnewline
\hline
\cel\corr{(iv) ESBMC} & \cel\corr{1.ª: ESBMC 8.1.0, Boolector 3.2.4, i7-13700K, Ubuntu 24.04. 2.ª: ESBMC 8.1.0--8.5.0, Boolector/Z3, máquina da autora, Ubuntu 24.04} & \cel\corr{1.ª: \texttt{--unwind 10 --no-unwinding-assertions --overflow-check}. 2.ª: \texttt{--unwind} por teste, com e sem \textit{unwinding assertions}, \texttt{--overflow-check --ub-shift-check}} & \cel\corr{limite $k$; sem \textit{unwinding assertions}, laços maiores que $k$ são truncados} \tabularnewline
\hline
\cel\corr{(v) Contraexemplos} & \cel\corr{1.ª: leitura do código. 2.ª: \textit{replay} com g++~13.3.0 e ASan/UBSan} & \cel\corr{\texttt{-std=c++17 -O0} \texttt{-fsanitize=address,undefined}} & \cel\corr{E/S simulada de forma determinística} \tabularnewline
\hline
\cel\corr{(vi) Reporte} & \cel\corr{GitHub (PX4-Autopilot)} & \cel\corr{---} & \cel\corr{---} \tabularnewline
\hline
\cel\corr{(vii) Validação cruzada} & \cel\corr{CBMC 5.95.1 (MiniSat)} & \cel\corr{opções equivalentes (Tabela~\ref{tab:equivalencia-opcoes})} & \cel\corr{mesmo \textit{harness} e mesmas pré-condições do ESBMC} \tabularnewline
\hline
\cel\corr{(viii) ArduPilot e CI} & \cel\corr{CBMC 5.95.1; GitHub Actions (\texttt{ubuntu-latest})} & \cel\corr{\texttt{--unwind 10 --signed-overflow-check --unsigned-overflow-check --bounds-check}} & \cel\corr{lógica equivalente no driver do ArduPilot} \tabularnewline
\hline
\end{tabular}
\end{table}

% =============================================================================
\section{Visão Geral da Abordagem}
\label{sec:visao-geral}
% =============================================================================

\corr{A Figura~\ref{fig:metodologia-visao-geral} sintetiza as oito etapas da metodologia e indica como elas foram aplicadas nas duas campanhas. O fluxo destaca a conferência com o código de produção, necessária para distinguir violações reproduzíveis de resultados decorrentes das limitações do modelo.}

%% LC 03 (resolvido 19/10/2026): seta "sem violacao" -> reporte removida.
%% [12/11/2026] Figura atualizada para englobar as duas campanhas e recolorida
%% para a paleta nude.
\begin{figure}[!tb]
\centering
\resizebox{0.92\linewidth}{!}{%
\begin{tikzpicture}[
  font=\linespread{1}\footnotesize,
  >={Stealth[length=2mm]},
  node distance=4mm,
  base/.style={draw=palEdge, rounded corners=2pt, align=center, inner sep=3pt,
    fill=palPaper},
  etapa/.style={base, text width=5.8cm, minimum height=7mm},
  firm/.style={etapa, fill=palFill},
  nova/.style={base, text width=4.8cm, minimum height=7mm, draw=palEdgeDark,
    fill=palStrong},
  dir/.style={base, text width=4.8cm, minimum height=7mm},
  conf/.style={base, text width=6.2cm, fill=palFill, draw=palEdgeDark, very thick},
  dec/.style={draw=palEdge, diamond, aspect=2.4, inner sep=1pt, fill=palPaper,
    align=center},
  desf/.style={base, text width=3.0cm, minimum height=10mm, fill=palStrong,
    draw=palEdgeDark},
  outn/.style={base, text width=3.0cm, minimum height=10mm, fill=palPaper},
  rot/.style={font=\scriptsize\itshape, align=center}]

  % --- espinha (i)-(v)
  \node[etapa] (e1) {(i) Seleção de módulos críticos};
  \node[firm, below=of e1] (e2) {(ii) Instrumentação com \textit{stubs} de \textit{hardware}};
  \node[firm, below=of e2] (e3) {(iii) Especificação de propriedades\\(\textit{datasheet}, documentação, código)};
  \node[etapa, below=of e3] (e4) {(iv) Configuração e execução do ESBMC};
  \node[etapa, below=of e4] (e5) {(v) Análise de contraexemplos};

  % --- conferencia compartilhada, com os dois modos
  \node[conf, below=6mm of e5] (conf) {\textbf{Conferência contra o código de produção}\\[1pt]
    1.\textsuperscript{a} campanha: leitura do código\\
    2.\textsuperscript{a} campanha: \textit{replay} do código real {\footnotesize+} ASan/UBSan};
  \node[dec, below=6mm of conf] (dec) {Reproduz em\\produção?};

  % --- desfechos
  \node[desf, below=10mm of dec] (def) {Defeito reproduzido / VAM};
  \node[outn, left=20mm of def] (fp) {Falso positivo (modelagem / ferramenta)};
  \node[outn, right=20mm of def] (ok) {Sem violação até o limite $k$};

  % --- coluna direita (vi)-(viii)
  \node[dir, right=2.4cm of e1.north east, anchor=north west] (e6) {(vi) Reporte à comunidade PX4\\(condicionado à revisão)};
  \node[nova, below=of e6] (e7) {(vii) Validação cruzada (CBMC)\\e verificação de \textit{patch}};
  \node[nova, below=of e7] (e8) {(viii) Reprodução (ArduPilot)\\e integração contínua};

  % --- setas da espinha e da coluna direita
  \foreach \a/\b in {e1/e2,e2/e3,e3/e4,e4/e5,e5/conf,conf/dec,e6/e7,e7/e8}
    \draw[->] (\a) -- (\b);

  % --- decisao -> desfechos
  \draw[->] (dec) -- node[rot, right] {sim} (def);
  \draw[->] (dec.west) -| node[rot, pos=0.25, above] {guarda\\omitida} (fp.north);
  \draw[->] (dec.east) -| node[rot, pos=0.25, above] {nenhuma\\violação} (ok.north);

  % --- defeito/VAM -> reporte (vi)
  \draw[->] (def.south) -- ++(0,-5mm) -| ([xshift=7mm]e6.east) -- (e6.east);

  % --- legenda
  \node[firm, text width=3mm, minimum height=3mm, below=12mm of fp.south west,
    anchor=north west] (l1) {};
  \node[right=1mm of l1, font=\scriptsize, anchor=west] (t1) {etapas de \textit{firmware}};
  \node[nova, text width=3mm, minimum height=3mm, right=5mm of t1] (l2) {};
  \node[right=1mm of l2, font=\scriptsize, anchor=west] (t2) {acréscimos deste trabalho};
  \node[conf, text width=3mm, minimum height=3mm, right=5mm of t2] (l3) {};
  \node[right=1mm of l3, font=\scriptsize, anchor=west] {conferência (1.\textsuperscript{a} e 2.\textsuperscript{a} campanhas)};
\end{tikzpicture}}
\caption{Metodologia de verificação em oito etapas, aplicada em duas campanhas.
As etapas~(i) e~(iv)--(vi) seguem o fluxo de Sousa et al.~\cite{Sousa2026}; as
etapas~(ii) e~(iii) são específicas de \textit{firmware} e as etapas~(vii)
e~(viii) são acréscimos deste trabalho. A conferência contra o código de
produção (etapa~v) é feita por leitura do código na primeira campanha e por
\textit{replay} do código real sob sanitizadores (ASan/UBSan) na segunda. Apenas
defeitos reproduzidos e violações alcançáveis no modelo (VAM) seguem para o
reporte~(vi), que depende de revisão; falsos positivos e resultados sem violação
encerram o fluxo do módulo. Os desfechos usam a escala do
Capítulo~\ref{cap:introducao}. Elaborado pela autora.}
\label{fig:metodologia-visao-geral}
\end{figure}

%% Prof. Lucas: "O que voce define como verificacao modular?" -- CORRIGIDO
\revisao{Foi adotada a verificação modular --- abordagem na qual o sistema é
decomposto em unidades funcionais (módulos) verificadas de forma independente, e
não como um único artefato monolítico. Cada módulo é isolado de suas dependências
por meio de abstrações (\textit{stubs}), e as propriedades de segurança são
verificadas no escopo local da unidade, o que reduz o tamanho do espaço de
estados analisado em cada execução do verificador~\cite{Cordeiro2012}.}

\revisao{A escolha decorre da dimensão da base de código.}
A árvore \texttt{src/} do PX4 soma aproximadamente 710 mil linhas de
código C/C++ (arquivos-fonte e cabeçalhos, contadas com \texttt{wc -l} no
commit 8be2590c15)\revisao{.
Essa dimensão torna a verificação monolítica impraticável devido à
explosão combinatória do espaço de estados~\cite{Biere1999}.}

%% Prof. Lucas: "O que voce define como remessa?" -- CORRIGIDO
\corr{A verificação da primeira campanha foi organizada em três rodadas. Neste
trabalho, uma rodada designa um conjunto de módulos verificados em um mesmo ciclo
de trabalho, agrupados por afinidade funcional ou por prioridade de criticidade.
A organização em rodadas permite escalonar o esforço de verificação, priorizando
primeiro os componentes cuja falha tem maior impacto sobre a segurança de voo e
ampliando a cobertura em ciclos subsequentes.}

%% LC 04 (resolvido 19/10/2026): definicao de harness em destaque.
\corr{\paragraph{Definição de \textit{harness}.} Neste trabalho, um
\textbf{\textit{harness} de verificação} é um arquivo C/C++ escrito
especificamente para a verificação, que (a)~declara as entradas do código sob
análise como valores não determinísticos, (b)~restringe essas entradas por
pré-condições (\textit{assumptions}), (c)~invoca as funções do módulo e
(d)~expressa as propriedades de segurança como \textit{assertions}. O
\textit{harness} tem uma função de entrada própria (\texttt{main()} na primeira
campanha; \texttt{harness\_main()} na segunda, porque o código original do PX4
já define uma \texttt{main} por meio de \texttt{ModuleBase}) e seleciona, por
uma variável de escolha, qual dos testes do arquivo será verificado.}

O ESBMC suporta verificação modular por contratos de função, nos
quais pré-condições e pós-condições são especificadas formalmente para cada
função\footnote{Documentação oficial do ESBMC sobre contratos de função: \url{https://esbmc.github.io/docs/function-contracts/}}. Neste trabalho, optou-se por uma abordagem
complementar: em vez de contratos explícitos, as propriedades foram expressas
diretamente como \textit{assertions} nos \textit{harnesses}.
%% LC 05 (resolvido 19/10/2026): justificativa reescrita por campanha.
\corr{Duas razões justificam essa escolha. Primeiro, as propriedades de
interesse vêm de documentos externos ao código (\textit{datasheets}, manuais e
especificações de protocolo) e se expressam naturalmente como condições sobre
o resultado de uma chamada, que é o que uma \textit{assertion} no
\textit{harness} captura. Segundo, um contrato precisaria ser anotado na própria
função verificada: na primeira campanha, as funções transcritas no
\textit{harness} são curtas e não justificariam a sobrecarga de contratos; na
segunda, anotar contratos exigiria modificar o código-fonte do PX4, o que essa
campanha justamente evita.}

\subsection{Etapa~(i) --- Seleção de Módulos Críticos}

A seleção dos módulos considerou dois critérios complementares: a relevância funcional para a segurança de voo e a viabilidade técnica da verificação isolada. O primeiro foi avaliado pela posição do módulo na cadeia de dados e de controle. Assim, receberam prioridade os drivers de sensores que fornecem dados ao estimador de estado (EKF2), os componentes de comunicação que processam comandos externos e os módulos de navegação responsáveis pela definição de trajetórias autônomas. Falhas nesses elementos podem afetar diretamente as decisões de voo. O segundo critério considerou a possibilidade de isolar o módulo das demais partes do \textit{firmware} por meio de abstrações, sem descaracterizar a lógica sob análise. \corr{Na segunda campanha, esses mesmos critérios orientaram a seleção de módulos de entrada, de sistema e de saída (Seção~\ref{sec:segunda-campanha}).}

%% LC 06 (resolvido 19/10/2026): "escrito manualmente" -> "assistencia de IA
%% generativa e revisadas pela autora".
%% [#2 RESOLVIDA 12/11/2026] A declaracao de IA e um \chapter* nao numerado
%% (pre-textual, sem \label); trocou-se o \ref por referencia textual.
\revisao{Cabe justificar a divisão em rodadas, já que o tempo total de execução do
verificador é da ordem de poucos minutos e, isoladamente, não a justificaria. O custo
dominante desta metodologia não é o tempo de máquina, e sim o \textbf{esforço humano} que o
antecede e o sucede: leitura do módulo, identificação das dependências a abstrair, construção
do \textit{harness}, derivação das propriedades a partir da documentação técnica e,
sobretudo, triagem dos contraexemplos.} \corr{Os \textit{harnesses} da primeira campanha somam
aproximadamente $2.800$ linhas, construídas com assistência de IA generativa e
revisadas pela autora, conforme a Declaração de Uso de Inteligência Artificial,
nas páginas iniciais desta dissertação, ao passo que a campanha completa executa
em cerca de dois minutos e meio.}

\revisao{É esse desequilíbrio que organiza o ciclo. Verificar em rodadas permite que o
aprendizado obtido em um conjunto de módulos --- quais abstrações são adequadas, quais
configurações do verificador cada perfil de módulo exige, quais padrões de contraexemplo são
recorrentes --- seja incorporado antes de investir esforço no conjunto seguinte. A necessidade
de \texttt{--unwind 300} no MAVLink e de \texttt{--no-library} no VTOL, por exemplo, só se
tornou conhecida durante a primeira rodada e informou a configuração das seguintes. A divisão
é, portanto, um mecanismo de organização do trabalho e de priorização por criticidade, e não
uma consequência do custo computacional.}

A organização em rodadas reflete essa priorização: a primeira rodada concentra
os drivers de sensores e comunicação de maior criticidade, cuja falha pode
comprometer toda a missão; as rodadas subsequentes ampliam a cobertura para
monitoramento, estimação de estado e navegação. Esse escalonamento permite
identificar \revisao{candidatos} de alto impacto antes de investir tempo de
verificação em módulos auxiliares.

\subsection{Etapa~(ii) --- Instrumentação com \textit{Stubs}}

%% LC 23 (resolvido 19/10/2026): etapa reescrita distinguindo as duas campanhas.
%% LC 22 (resolvido 19/10/2026): explicacao de nao determinismo mantida so aqui.
\corr{A instrumentação tem por finalidade abstrair as entradas e saídas de \textit{hardware}, as chamadas ao sistema operacional NuttX e as dependências de outros módulos. Para isso, utiliza-se a primitiva \texttt{nondet\_T()}, em que \texttt{T} representa o tipo do dado. Essa primitiva permite ao verificador considerar qualquer valor representável por esse tipo, incluindo condições extremas, dados inválidos e respostas inesperadas do \textit{hardware}, sem exigir a enumeração manual de casos de teste~\cite{Menezes2024}. Por exemplo, \texttt{nondet\_uint8()} representa qualquer valor entre $0$ e $255$.}

\corr{As duas campanhas aplicam essa abstração de formas diferentes. Na
\textbf{primeira}, os \textit{harnesses} não contêm funções de \textit{stub}: os
bytes que o driver leria do barramento (SPI, I2C ou UART) são declarados como
valores não determinísticos na própria função de teste e passados diretamente ao
trecho do driver transcrito no \textit{harness}
(Listagens~\ref{lst:stub_i2c} e~\ref{lst:main_imu}). Na \textbf{segunda}, como o
código verificado é o original, as funções de que ele depende precisam existir:
foram escritos \textit{stubs} de fato --- \texttt{read}, \texttt{write} e
\texttt{select} da UART, o barramento I2C, o \textit{middleware} uORB, o tempo, os
parâmetros e o registro de mensagens ---, que preservam a assinatura original e
devolvem valores não determinísticos (Seção~\ref{sec:segunda-campanha}).}

%% LC 07 (resolvido 19/10/2026): pre-condicoes tabuladas e justificadas.
\corr{Os valores não determinísticos, porém, não ficam irrestritos: cada
\textit{harness} restringe parte das entradas por pré-condições
(\texttt{\_\_ESBMC\_assume}), seja para refletir a faixa que o código de produção
de fato recebe, seja para concentrar a busca numa região de interesse. Essas
restrições delimitam o alcance de cada veredito e estão tabuladas, com a
justificativa de cada uma, na Tabela~\ref{tab:hipoteses-harness}. Pré-condições
estreitas demais podem mascarar defeitos; pré-condições largas demais admitem
estados que o driver real rejeita e produzem falsos positivos --- o mecanismo do
candidato do contador de FIFO. Por isso, toda violação passa pela conferência
contra o código de produção da etapa~(v).}

\begin{table}[!htb]
\centering
\caption{\corr{Principais pré-condições dos \textit{harnesses} da primeira campanha e sua justificativa}}
\label{tab:hipoteses-harness}
\footnotesize
\begin{tabular}{|p{2.3cm}|p{5.6cm}|p{6.0cm}|}
\hline
\textbf{Módulo} & \textbf{Pré-condição (\texttt{\_\_ESBMC\_assume})} & \textbf{Justificativa} \\
\hline
\corr{GPS uBlox} & \corr{\texttt{0 < input\_len <= GPS\_DUMP\_DATA\_SIZE + 10}; \texttt{dump\_buffer.len >= GPS\_DUMP\_DATA\_SIZE - 5}} & \corr{tamanho do vetor de entrada do \textit{harness}; estado inicial próximo da fronteira (discutido na Seção~\ref{sec:exemplo-ilustrativo})} \\
\hline
\corr{IMU BMI088} & \corr{eixos do acelerômetro e do giroscópio em \texttt{[INT16\_MIN, INT16\_MAX]}; bytes de temperatura e FIFO irrestritos} & \corr{faixa do tipo; não restringe nada além do tipo} \\
\hline
\corr{MAVLink} & \corr{\texttt{0 < packet\_len <= BUFFER\_SIZE + 50}} & \corr{admite, de propósito, pacotes maiores que o \textit{buffer} (280 bytes), para exercitar a guarda de capacidade} \\
\hline
\corr{Expo RC} & \corr{entrada em $[-1;\,1]$ e expoente em $[0;\,1]$, finitos} & \corr{faixa normalizada do controle remoto e do parâmetro de expo} \\
\hline
\corr{RC Update} & \corr{\texttt{min} em $[800;1200]$, \texttt{trim} em $[1400;1600]$, \texttt{max} em $[1800;2200]$, \texttt{min < trim < max}} & \corr{faixas de calibração de canal em microssegundos} \\
\hline
\corr{PWM Output} & \corr{saída em $[\texttt{PWM\_MOTOR\_MIN}; \texttt{PWM\_MOTOR\_MAX}]$ no teste nominal; $[-10^4;\,10^4]$ no teste de robustez} & \corr{faixa de motor; teste de robustez com entrada fora da faixa} \\
\hline
\corr{EKF2} & \corr{variâncias finitas, não negativas e limitadas por \texttt{VAR\_VEL\_MAX}; \texttt{min\_var <= max\_var}} & \corr{limites das variâncias usados no estimador} \\
\hline
\corr{Magnetômetro} & \corr{variâncias finitas, em $[0;\,\texttt{MAX\_VAR\_ALLOWED}]$, com soma positiva; ângulos em $[0;\,\pi]$} & \corr{evita divisão por zero no ganho; faixa de ângulo} \\
\hline
\corr{Navigator} & \corr{índices \texttt{<= MAX\_MISSION\_ITEMS}; distância em $[0;\,10^5]$ m; raio de aceitação em $[0{,}1;\,1000]$ m} & \corr{tamanho máximo de missão; faixas físicas plausíveis} \\
\hline
\corr{VTOL} & \corr{densidade do ar em $[0{,}1;\,3]$; pesos em $[0{,}1;\,100]$; velocidade de transição em $[5;\,50]$ m/s} & \corr{faixas físicas plausíveis} \\
\hline
\corr{Barômetro} & \corr{pressão bruta em $[50000;\,110000]$ Pa; \textit{offsets} limitados; prioridade em $[-200;\,200]$} & \corr{faixa de operação de barômetros; teste de prioridade além da faixa válida} \\
\hline
\corr{Battery Monitor} & \corr{SoC em $[0;\,1]$; limiares ordenados; tensão de célula em $[3{,}0;\,4{,}5]$ V; 1 a 12 células} & \corr{faixas de bateria LiPo e de configuração} \\
\hline
\end{tabular}
\end{table}

%% LC 08 (resolvido 19/10/2026): citacao Godefroid2005 removida.
\corr{Para cada módulo, um arquivo de \textit{harness} centraliza as entradas
não determinísticas, as pré-condições, as funções de teste e a função de
entrada.}

\subsection{Etapa~(iii) --- Especificação de Propriedades}

As propriedades verificadas foram organizadas em duas categorias: automáticas e customizadas. As propriedades automáticas são habilitadas por opções de linha de comando do ESBMC e abrangem classes gerais de defeitos, como acessos fora dos limites de vetores e \textit{buffers}, desreferências de ponteiros nulos e \textit{overflows} aritméticos em inteiros~\cite{Cordeiro2012}. Essas verificações dispensam \textit{assertions} específicas para cada módulo e podem ser aplicadas de maneira uniforme, desde que as opções correspondentes estejam habilitadas.

%% LC 09 (resolvido 19/10/2026): rastreabilidade atenuada (1.a) e afirmada (2.a).
As propriedades customizadas capturam requisitos específicos de cada módulo,
derivados de datasheets de fabricantes e da especificação do protocolo PX4. São
expressas como assertions inseridas no harness, imediatamente após a chamada às
funções verificadas. \corr{A rastreabilidade entre cada \textit{assertion} e o
requisito de origem é registrada na Tabela~\ref{tab:proveniencia}. Nos
\textit{harnesses} da primeira campanha, os comentários junto às funções de
teste descrevem a especificação em linguagem natural (por exemplo,
``ESPECIFICAÇÃO: Converter dados raw de temperatura para Celsius''), mas não
indicam a seção nem a página do documento. Na segunda campanha, cada
\textit{harness} anota, no cabeçalho, a origem de cada pré-condição e de cada
propriedade, com seção e página do documento quando ele existe (por exemplo,
o manual da RoboClaw, Rev.~6.0, pp.~48--49).}

\subsection{Etapa~(iv) --- Configuração e Execução do ESBMC}

A configuração do ESBMC depende, entre outros fatores, do limite de desenrolamento e do \textit{solver} SMT selecionado. O parâmetro \texttt{--unwind} determina quantas iterações de cada laço serão exploradas. Limites menores reduzem o custo da análise, mas podem deixar de alcançar defeitos que ocorram em iterações posteriores. Já a opção \texttt{--default-solver} seleciona o \textit{solver} utilizado e pode afetar tanto o tempo de execução quanto a obtenção de um resultado dentro do tempo-limite. Nesta pesquisa, o Z3 não produziu um veredito dentro do \textit{timeout} nos testes indicados para GPS, MAVLink e EKF2 com \texttt{--unwind 10}; por esse motivo, o Boolector foi adotado nas execuções correspondentes da primeira campanha (Seção~\ref{sec:ameacas-validade}).
%% LC 10 (resolvido 19/10/2026): bit-blasting de ponto flutuante do ESBMC.
\corr{Como o Boolector não suporta a teoria de ponto flutuante do padrão
SMT-LIB, o ESBMC codifica as operações IEEE-754 de precisão simples e dupla como
operações sobre vetores de bits (\textit{bit-blasting} de ponto flutuante),
preservando a semântica de arredondamento, os valores especiais (NaN, infinitos)
e os números subnormais; com solvers que suportam essa teoria, como o Z3, a
codificação pode usá-la diretamente~\cite{Gadelha2020}. Essa codificação
precisa é essencial para os módulos de sensores e de estimação, que usam
aritmética de ponto flutuante.}

A configuração é iterativa: parte-se de parâmetros padrão e ajusta-se conforme
os resultados. Uma violação de asserção de desenrolamento (\textit{unwinding
assertion}) indica que o limite atual é insuficiente para cobrir os laços do
módulo, exigindo aumento do parâmetro \texttt{--unwind}, como ocorreu
no MAVLink ($k = 10 \rightarrow 300$). Quando as funções da biblioteca matemática
tornam o modelo intratável, a opção \texttt{--no-library} desativa a biblioteca C
abstrata do ESBMC: funções como \texttt{sinf} e \texttt{sqrtf} passam a
retornar valores não determinísticos, uma sobreaproximação que não deixa escapar
violações das demais verificações, mas deixa de verificar as propriedades
internas à biblioteca --- caso do VTOL~\cite{Menezes2024}.

\subsection{Etapa~(v) --- Análise de Contraexemplos e Conferência contra o Código de Produção}

Quando o ESBMC reporta uma violação, o contraexemplo gerado
apresenta valores concretos de entrada que conduzem ao estado de falha.
\revisao{A análise começa pela leitura do traço: identifica-se qual propriedade foi
violada, em que linha, e qual atribuição de valores conduziu ao estado. Nessa
etapa o contraexemplo ainda é um \textbf{candidato}, na acepção fixada no
Capítulo~\ref{cap:introducao} --- não há, até aqui, evidência de que o estado
seja alcançável fora do ambiente de verificação.}

\revisao{A classificação e a avaliação de impacto ocorrem depois, com base nas
consequências da falha no contexto de voo. Um \textit{overflow} aritmético em
variável de controle de atitude tem impacto distinto de uma violação de asserção
em caminho inalcançável em condições normais. Essa triagem guia quais resultados
seguem para as etapas seguintes e com qual prioridade.}

\paragraph{Conferência contra o código de produção.}
\label{subsec:etapa-replay}

\revisao{Esta conferência decide se um candidato é violação alcançável no modelo ou falso
positivo, e é a que sustenta a distinção que organiza os resultados deste
trabalho. Um contraexemplo é sempre válido \emph{para o modelo codificado}; o que
se verifica é se o estado violado persiste quando o
\textit{harness} sai de cena e o código de produção entra em seu lugar.}
Ela pode ser feita de duas formas: por \textit{replay} ou por leitura do
código de produção.

%% LC 11 (resolvido 19/10/2026): protocolo reescrito pelo que foi executado na
%% segunda campanha (candidatos extras, RoboClaw e batmon).
\corr{No \textit{replay}, o arquivo-fonte \emph{original} do driver é compilado
com um compilador comum --- g++~13.3.0, com \texttt{-std=c++17 -O0} --- junto com
\textit{stubs} determinísticos que substituem o \textit{hardware} e injetam os
valores exatos do contraexemplo: a UART simulada registra cada pacote enviado e
devolve as respostas escolhidas pelo programa de \textit{replay}. O programa é
executado, sob os sanitizadores de endereço (ASan) e de comportamento indefinido
(UBSan), e registra o efeito observável do caminho do contraexemplo, como o byte
efetivamente enviado ao atuador ou o acesso a um índice fora dos limites.}

\corr{O critério de decisão é binário. Se a execução do código original
reproduz o efeito previsto pelo contraexemplo --- por exemplo, um byte fora da
faixa especificada no pacote enviado ---, o candidato é classificado como
\textbf{violação alcançável no modelo} e, conforme a análise de alcançabilidade
no uso real, como defeito. Se não reproduz, investiga-se a razão: se o código de
produção contém guarda que o \textit{harness} omitiu, o candidato é
\textbf{falso positivo}. Duas verificações de sanidade acompanham o
\textit{replay}: confirmar que a execução alcança a função analisada, e não
termina antes, no código de preparação; e confirmar que o cenário crítico é
derivado da entrada injetada, e não fixado por atribuição direta no programa de
\textit{replay} --- caso em que se estaria demonstrando aritmética, e não
alcançabilidade.}

%% [#3 RESOLVIDA 12/11/2026] O replay de FIFO/GPS nao foi executado; a diferenca
%% de procedimento da primeira campanha fica registrada como limitacao.
Nos três contraexemplos da primeira campanha, a conferência foi
feita por leitura do código de produção: o candidato do contador de FIFO foi
refutado pela inspeção do driver do BMI088, e a inalcançabilidade do estado
violador do GPS veio da revisão do mantenedor e da análise do código-fonte,
depois verificada por prova indutiva (Seção~\ref{subsec:gps-inducao}).
\corr{Na segunda campanha, o \textit{replay} com sanitizadores foi executado para
os candidatos alcançáveis dos oito módulos adicionais, do driver da RoboClaw e do
batmon (Seção~\ref{sec:segunda-campanha}); os programas de \textit{replay} e seus
registros estão no repositório do trabalho (\texttt{candidatos\_extras\_v5/},
\texttt{testes\_hardware/roboclaw/} e \texttt{novo\_achado\_batmon/}). Registra-se,
como limitação, que o \textit{replay} executável não foi aplicado aos candidatos
do contador de FIFO e do GPS da primeira campanha, cuja conferência se deu apenas
por leitura do código de produção; essa diferença de procedimento é retomada nas
ameaças à validade (Seção~\ref{sec:ameacas-validade}).}

\subsection{Etapa~(vi) --- Reporte \corr{à} Comunidade PX4}

%% LC 12 (resolvido 19/10/2026): desfecho de cada reporte e momento do reporte.
\corr{Os candidatos que sobrevivem à conferência são reportados como
\textit{issues} no repositório oficial do PX4 no GitHub, acompanhados de:
descrição do achado, contraexemplo com valores concretos, trecho de código
afetado, referência ao documento que define o comportamento esperado e, quando
houver, proposta de correção. Esse formato permite aos mantenedores reproduzir o
problema sem acesso ao ambiente de verificação~\cite{PX4_GitHub}.}

\corr{O reporte não garante aceitação --- e, nesta pesquisa, os da primeira
campanha não foram aceitos como defeitos. A observação de temperatura do BMI088
(\textit{Issue}~\#26865) foi recusada pelos mantenedores, que não reconheceram a
faixa operacional do \textit{datasheet} como restrição de software
(Seção~\ref{subsec:imu-temperatura}); o PR~\#27575, que tratava da temperatura e
do FIFO, também não foi aceito; e a guarda proposta para o GPS
(\textit{Issue}~\#26866, PR~\#27576) foi enquadrada pelo mantenedor como
endurecimento defensivo, depois que a revisão mostrou o estado violador
inalcançável no driver de produção. A revisão pela comunidade é, assim, a única
etapa capaz de elevar um achado à condição de defeito confirmado, e também a que
pode recusá-lo.}

\corr{O momento do reporte pesou nesses desfechos. Para o GPS e para o contador
de FIFO, a \textit{issue} foi aberta antes da conferência contra o código de
produção, pedindo a própria triagem; a conferência veio depois, e no caso do FIFO
rebaixou o candidato a falso positivo. A lição incorporada à segunda campanha foi
reportar apenas depois da conferência, do \textit{replay}, da consulta à
especificação com página citada e da verificação de que o código ainda é o mesmo
na branch principal e de que não há \textit{issue} aberta sobre o assunto. Em
coerência com essa lição e com a orientação da pesquisa, nenhum achado da segunda
campanha foi reportado sem a prévia revisão do orientador. No caso do RoboClaw,
em que o driver havia sido reescrito na branch principal, o relato indicou
explicitamente o ramo de \textit{release} em que o código analisado permanece
(Seção~\ref{subsec:2c-roboclaw}).}

\subsection{Etapa~(vii) --- \revisao{Validação Cruzada com Verificador Independente} e Verificação de \textit{Patch}}
\label{subsec:etapa-validacao-cruzada}

\revisao{Os candidatos são submetidos a um segundo
verificador, o CBMC~5.95.1~\cite{Kroening2014}, sobre o mesmo módulo e com opções
equivalentes. O objetivo é distinguir defeito de artefato de \textit{toolchain}:
se apenas um dos verificadores reporta a violação, a hipótese de artefato ganha
força; se ambos reportam, ela perde.}

\revisao{Cabe uma ressalva de escopo, aplicada de forma sistemática no
Capítulo~\ref{cap:experimentos}: a concordância entre dois verificadores sobre o
mesmo \textit{harness} corrobora a codificação do problema, e não a
alcançabilidade do estado no driver de produção. Os dois partem da mesma
abstração e das mesmas pré-condições. Acresce que os verificadores podem
classificar o mesmo defeito sob propriedades formais distintas, conforme tratem a
conversão entre tipos com e sem sinal --- o que é informação adicional, e não
contradição.} \corr{Na segunda campanha, a comparação foi estendida a todos os
\textit{harnesses}, e não apenas aos candidatos, para medir também a
aplicabilidade de cada ferramenta ao código do PX4
(Seção~\ref{subsec:2c-protocolo}).}

Quando uma correção é proposta, o mesmo \textit{harness} é reverificado com
o \textit{patch} aplicado, nas mesmas pré-condições: sem o \textit{patch}, a
violação deve aparecer; com ele, deve desaparecer. A comparação só isola o efeito
da guarda se as pré-condições forem as mesmas nos dois casos
(Seção~\ref{subsec:patch-gps}).

%% LC 13 (resolvido 19/10/2026): etapa renomeada; reproducao so com CBMC.
\subsection{Etapa~(viii) --- \corr{Reprodução em Outro Autopiloto e Integração Contínua}}
\label{subsec:etapa-ardupilot-ci}

Quando a lógica de um achado tem equivalente em outro autopiloto, ele é
conferido também lá. \corr{Na primeira campanha, essa reprodução foi feita
apenas com o CBMC~5.95.1, e não com o ESBMC; trata-se, portanto, de uma
reprodução da observação em outro código, e não de uma verificação completa do
outro autopiloto.} A observação de temperatura do BMI088 foi reproduzida
no driver equivalente do ArduPilot~\cite{ArduPilot}
(\texttt{AP\_InertialSensor\_BMI088.cpp}, linhas 383--385; \textit{Issue}
\#33272), onde um mantenedor encerrou a \textit{issue} citando a faixa de
medição de $-104$ a $150\,^{\circ}$C do \textit{datasheet}. A reprodução serve
apenas como verificação de consistência: os dois projetos recusaram a
observação.
%% [#5 RESOLVIDA 12/11/2026] A repeticao com o ESBMC 8.5 era opcional; declarada
%% como refinamento futuro.
\corr{A repetição dessa reprodução com o ESBMC~8.5, para que a etapa empregasse
as duas ferramentas também no ArduPilot, fica como refinamento futuro; como
executada, a reprodução usou somente o CBMC.}

%% LC 14 (resolvido 12/11/2026): workflow lido no fork (esbmc-verify.yml).
%% Observacao de fidelidade: o workflow instala e roda o CBMC (nao o ESBMC) sobre
%% versoes reduzidas dos harnesses do IMU e do GPS -- descrito como smoke check.
%% So o tempo medido de uma rodada permanece pendente (aba Actions).
\corr{Por fim, um \textit{workflow} do GitHub Actions
(\texttt{esbmc-verify.yml}), em um \textit{fork} do PX4
(\url{https://github.com/marcela2881/PX4-Autopilot}), executa a cada
\textit{push} e \textit{pull request} para a \textit{branch} \texttt{main} uma
verificação de fumaça (\textit{smoke check}) com o CBMC, em \textit{runner}
\texttt{ubuntu-latest}, sobre versões reduzidas dos \textit{harnesses} do IMU
BMI088 e do GPS uBlox (\texttt{--unwind 10}, com checagem de \textit{overflow}
com e sem sinal e de limites). O \textit{workflow} reexecuta essa verificação a
cada alteração do código, servindo de salvaguarda contra regressões; emprega o
CBMC, e não o ESBMC, por ser mais simples de instalar no ambiente de integração
contínua.}

% =============================================================================
\section{Exemplo Ilustrativo: GPS uBlox}
\label{sec:exemplo-ilustrativo}
% =============================================================================

A Listagem~\ref{lst:exemplo-harness-gps} apresenta o fragmento central do harness construído para o GPS uBlox, destacando os elementos essenciais da metodologia: a declaração da entrada não determinística (\revisao{\texttt{nondet\_size\_t()}}), a \textit{assumption} que restringe o domínio (\texttt{\_\_ESBMC\_assume}) e a \textit{assertion} que expressa a propriedade de segurança verificada. O código completo de todos os harnesses e os logs de execução do ESBMC estão disponíveis em \url{https://github.com/marcela2881/dissertacao-px4-esbmc}.

\begin{lstlisting}[
  caption={\revisao{Fragmento do harness de verificação do GPS uBlox, função \texttt{test\_gps\_real\_bit\_operation}, transcrito de \texttt{test\_gps\_driver.cpp}.}},
  label={lst:exemplo-harness-gps},
  numbers=left,
  frame=single,
  language=C]
// Entrada simbolica - o solver escolhe o valor
size_t input_len = nondet_size_t();

// Restricao do dominio: o buffer de entrada tem
// GPS_DUMP_DATA_SIZE + 10 = 210 bytes
__ESBMC_assume(input_len > 0 && input_len <= GPS_DUMP_DATA_SIZE + 10);

uint8_t input_data[GPS_DUMP_DATA_SIZE + 10];
gps_dump_s dump_buffer;

// Estado inicial do buffer, proximo da fronteira
dump_buffer.len = nondet_uint8();
__ESBMC_assume(dump_buffer.len >= GPS_DUMP_DATA_SIZE - 5);

// Execucao da funcao real do PX4 sobre a entrada simbolica
dumpGpsData(input_data, input_len, gps_dump_comm_mode_t::Full,
            /* msg_to_gps_device = */ true,
            &dump_buffer, gps_dump_comm_mode_t::Full);

// Propriedade de seguranca: o comprimento acumulado nao deve
// ultrapassar GPS_DUMP_DATA_SIZE, considerado o bit 7 de sinalizacao
assert(dump_buffer.len == 0 ||
       dump_buffer.len <= GPS_DUMP_DATA_SIZE ||
       (dump_buffer.len & 0x7F) <= GPS_DUMP_DATA_SIZE);
\end{lstlisting}

\revisao{As duas pré-condições delimitam o espaço de entradas e, com ele, o alcance
do veredito. A primeira restringe \texttt{input\_len} ao tamanho do vetor de
entrada declarado no próprio \textit{harness}, $210$ bytes, evitando que o
verificador reporte acessos fora dos limites do vetor de teste em vez do defeito
sob investigação. A segunda posiciona \texttt{dump\_buffer.len} a no máximo cinco
unidades da fronteira de $200$ bytes, concentrando a busca na vizinhança em que a
subtração sem sinal pode dar a volta. Trata-se de decisão de modelagem que
restringe a busca: violações que exigissem estados iniciais distantes da
fronteira não seriam alcançadas por este \textit{harness}.}

%% LC 15 (resolvido 19/10/2026): tautologia declarada; propriedade de memoria do
%% memcpy e a verificada.
\corr{A \textit{assertion} final da Listagem~\ref{lst:exemplo-harness-gps} é,
na prática, uma tautologia: o terceiro termo, \texttt{(dump\_buffer.len \& 0x7F)
<= GPS\_DUMP\_DATA\_SIZE}, é sempre verdadeiro, pois o valor mascarado com
\texttt{0x7F} nunca passa de $127$, que é menor que $200$. A \textit{assertion}
não detecta, portanto, nenhuma violação. A violação reportada pelo ESBMC vem da
verificação automática de limites dentro do \texttt{memcpy}
(\textit{dereference failure}: acesso fora dos limites do objeto), que é a
propriedade efetivamente verificada neste teste. Uma propriedade de saída
significativa seria \texttt{dump\_buffer.len <= GPS\_DUMP\_DATA\_SIZE} após a
chamada, sem os termos disjuntivos.}

%% [#4 RESOLVIDA 12/11/2026] Harness reexecutado: ESBMC 8.2.0, Z3, ~389 s,
%% VERIFICATION SUCCESSFUL. Log em resultados/GPS/gps_prop_limpa.txt.
\corr{Para separar o que é defeito do que é artefato dessa modelagem, o
\textit{harness} foi reexecutado com a propriedade de saída na forma não
tautológica (\texttt{dump\_buffer.len <= GPS\_DUMP\_DATA\_SIZE}) e partindo de um
estado inicial válido (\texttt{dump\_buffer.len < GPS\_DUMP\_DATA\_SIZE}). Nessas
condições o ESBMC~8.2.0 (solver Z3) verifica \textsc{verification successful} em
cerca de $389$~s: a lógica transcrita do \texttt{dumpGpsData} respeita o limite a
partir de qualquer estado válido. A violação observada na versão exploratória
decorre, portanto, exclusivamente do estado inicial acima da capacidade admitido
pela pré-condição --- estado inalcançável no driver de produção, como confirma a
prova indutiva da Seção~\ref{subsec:gps-inducao}. Os dois resultados são
consistentes: o modelo só viola a propriedade quando parte de um estado que o
código de produção nunca alcança, o que sustenta a classificação do achado como
violação alcançável no modelo, e não como defeito.}

%% LC 16 (resolvido 19/10/2026): violacao e admitida pela pre-condicao.
\corr{A segunda pré-condição merece uma ressalva adicional. Como ela admite
\texttt{dump\_buffer.len} entre $195$ e $255$, o \textit{harness} já parte de
estados em que o preenchimento está \emph{acima} da capacidade de $200$ bytes ---
o contraexemplo usa $207$. A violação, portanto, não é derivada a partir de um
estado inicial válido: ela é admitida pela pré-condição. O que o contraexemplo
mostra é que, \emph{se} o preenchimento passar da capacidade, a guarda do trecho
transcrito falha; não mostra que esse estado é alcançável. Foi exatamente essa
pergunta que a revisão do mantenedor e a prova indutiva da
Seção~\ref{subsec:gps-inducao} responderam negativamente para o driver de
produção.}

%% LC 17 (resolvido 19/10/2026): retirado "fiel"; duas diferencas declaradas.
\corr{O \textit{harness} também difere do driver em dois pontos, que impedem
chamá-lo de transcrição fiel. Primeiro, a capacidade: $200$ bytes no
\textit{harness} (\texttt{GPS\_DUMP\_DATA\_SIZE}) e $79$ no PX4
(\texttt{sizeof(dump\_data->data)}, em \texttt{msg/GpsDump.msg}). Segundo, os
tipos dos operandos: no \textit{harness}, a subtração é feita em \texttt{int}
(macro \texttt{int} menos \texttt{uint8\_t} promovido), e o resultado negativo é
convertido para \texttt{size\_t}; no PX4, a subtração já é em \texttt{size\_t}.
O mecanismo do \textit{underflow} é análogo nos dois, mas os valores e o
caminho não são idênticos; por isso, o resultado vale para o modelo, e a
conclusão sobre o driver depende da conferência da etapa~(v). Essa limitação,
comum a toda a primeira campanha, motivou a segunda campanha, em que o código
verificado é o original (Seção~\ref{sec:segunda-campanha}).}

%% [RESOLVIDO] Os tres numeros corrigidos (1.712 LoC; 64,9 s; oito etapas).

Para concretizar a aplicação da metodologia, esta seção percorre
\revisao{as oito etapas} tendo o driver GPS uBlox como objeto de verificação. O driver,
com \revisao{$1.712$ linhas de código} (o \textit{harness}
correspondente tem 294 linhas), é responsável pelo parsing de mensagens
do protocolo UBX recebidas via UART e pela atualização das estruturas de dados de
posição e velocidade consumidas pelo estimador EKF2.

\textbf{Seleção.} O GPS uBlox foi incluído na primeira rodada por sua posição
crítica na cadeia de navegação: coordenadas corrompidas podem comprometer o
Return-to-Launch e outras funções de failsafe que dependem de posição
absoluta~\cite{Wanner2024}.

\textbf{Instrumentação.} \corr{Os bytes que o driver receberia da UART foram
declarados como valores não determinísticos na função de teste e passados ao
trecho transcrito do driver}, simulando tanto pacotes UBX bem formados
quanto dados corrompidos ou truncados. A função \texttt{dumpGpsData()}, que copia
pacotes recebidos para um buffer de diagnóstico, foi incluída no harness por
conter aritmética sobre tipos sem sinal em sua verificação de espaço disponível.

\textbf{Propriedades.} Além das verificações automáticas de segurança de
memória, foi especificada a propriedade~\eqref{eq:prop_gps}: o comprimento do
buffer de diagnóstico não deve superar \texttt{GPS\_DUMP\_DATA\_SIZE} (200
bytes no \textit{harness}; no driver de produção, o registro tem 79 bytes)
após qualquer operação de escrita.

\textbf{Execução.} A verificação foi executada com Boolector 3.2.4,
\texttt{--unwind 10} e \texttt{--overflow-check}, \revisao{concluindo em $64{,}9$
segundos, como mediana de três execuções no servidor descrito na
Seção~\ref{sec:config-experimental}}.

\textbf{Análise.} \revisao{O ESBMC parou na primeira violação --- um
underflow aritmético: quando \texttt{dump\_data->len} supera 200, a expressão
\texttt{GPS\_DUMP\_DATA\_SIZE - dump\_data->len}, em que
\texttt{GPS\_DUMP\_DATA\_SIZE} é um \texttt{\#define} de valor 200 (do tipo
\texttt{int}) e \texttt{len}, um \texttt{uint8\_t} promovido a \texttt{int},
produz um valor negativo ($200 - 207 = -7$ no contraexemplo), que é convertido
para \texttt{size\_t} --- um valor enorme --- na comparação com
\texttt{write\_len}, tornando a guarda
de proteção sempre falsa. O \texttt{memcpy} subsequente copia dados sem
restrição de limite, violando os limites do objeto alocado. No driver
de produção (\texttt{gps.cpp}, linha 741), a expressão correspondente é
\texttt{sizeof(dump\_data->data) - dump\_data->len}, avaliada em
\texttt{size\_t} (sem sinal), sobre um buffer de 79 bytes
(\texttt{msg/GpsDump.msg}).}
\revisao{O contraexemplo, produzido pela função
\texttt{test\_gps\_real\_bit\_operation} (\texttt{test\_choice = 3}), tem
\texttt{input\_len = 10} e \texttt{dump\_buffer.len = 207}. A revisão do
mantenedor do PX4 no PR~\#27576 e a análise do código de produção mostraram
que o estado violador não ocorre no driver real: o laço de escrita limita
\texttt{write\_len} ao espaço livre e zera \texttt{len} após publicar o
registro. Essa inalcançabilidade foi depois verificada formalmente por prova
indutiva (Seção~\ref{subsec:gps-inducao}).}

\textbf{Reporte.} O \revisao{achado} foi reportado à comunidade PX4 via Issue
\#26866\footnote{\url{https://github.com/PX4/PX4-Autopilot/issues/26866}},
acompanhado do contraexemplo completo, e a correção foi proposta como uma
guarda no PR~\#27576, que o mantenedor enquadrou como
endurecimento defensivo~\cite{PX4_GitHub}.

\textbf{Validação cruzada e \textit{patch}.} O CBMC~5.95.1 sinaliza a
mesma violação apenas com \texttt{--unwind} 2, 3 e 5; com \texttt{--unwind 10},
não terminou em 30 minutos. Com as mesmas pré-condições, o ESBMC e o CBMC
acusam a violação sem o \textit{patch} e não a acusam com o \textit{patch} do PR
nem com a forma sugerida pelo mantenedor (Seção~\ref{subsec:patch-gps}).
\corr{Na comparação refeita em 19/10/2026 com o ESBMC~8.5.0 e o tempo-limite de
$600$~s, o ESBMC voltou a encontrar a violação em cerca de $137$~s (mediana de
três execuções) e o CBMC não terminou dentro do limite
(Seção~\ref{subsec:2c-protocolo}).}

\textbf{Integração contínua.} O \textit{harness} do GPS roda no
\textit{workflow} de integração contínua do \textit{fork} a cada \textit{push}
(Seção~\ref{subsec:etapa-ardupilot-ci}). A etapa de reprodução no ArduPilot não
se aplica a este achado, já que ela foi feita apenas para a observação de
temperatura do BMI088.

% =============================================================================
\section{Seleção dos Módulos para Verificação}
\label{sec:selecao-modulos}
% =============================================================================

\revisao{O procedimento de seleção combinou duas etapas de natureza distinta, e convém explicitá-las separadamente. A primeira é o \textbf{levantamento do conjunto candidato}, que não foi algorítmico: os módulos foram indicados pela equipe do projeto SAGRES em reunião, com base na experiência operacional da equipe sobre quais componentes são mais determinantes para missões de sensoriamento remoto em áreas remotas. Trata-se, portanto, de elicitação junto a especialistas do domínio, e não de um estudo sistemático de criticidade --- distinção que se registra aqui para que o leitor possa avaliar o alcance da amostra (Seção~\ref{sec:ameacas-validade}, no Capítulo~\ref{cap:experimentos}).}

\revisao{A segunda etapa é a \textbf{filtragem por critérios objetivos}, essa sim sistematizável e reproduzível, formalizada no Algoritmo~\ref{alg:selecao}.} \revisao{Cada candidato indicado foi submetido pela autora a três verificações sequenciais: leitura do código-fonte, para avaliar a complexidade interna e a superfície de entrada não confiável; leitura da documentação e do datasheet, para identificar requisitos de segurança explícitos; e avaliação da viabilidade técnica de isolamento por meio de \textit{stubs}. Apenas os módulos que atendem aos três critérios são incluídos no conjunto final.}

\begin{algorithm}[H]
\caption{Procedimento de seleção de módulos críticos do PX4}
\label{alg:selecao}
\begin{algorithmic}[1]
\Require Repositório PX4-Autopilot (ramo main, commit 8be2590c15), datasheets dos componentes, indicações do coordenador SAGRES
\Ensure Conjunto $M$ de módulos selecionados para verificação
\State $M \gets \emptyset$
\State $\text{candidatos} \gets$ módulos indicados pelo coordenador SAGRES
\ForAll{$m \in \text{candidatos}$}
    \State ler código-fonte de $m$ \Comment{análise da lógica interna}
    \State $\text{superficie}(m) \gets$ pontos de entrada de dado não confiável em $m$
    \State ler datasheet e documentação associados a $m$
    \State $\text{requisitos}(m) \gets$ requisitos de segurança extraídos dos documentos
    \State $\text{isolavel}(m) \gets$ avaliação da viabilidade de abstração via \textit{stubs}
    \If{$\text{superficie}(m) \neq \emptyset \,\land\, \text{requisitos}(m) \neq \emptyset \,\land\, \text{isolavel}(m) = \text{verdadeiro}$}
        \State $M \gets M \cup \{m\}$
    \EndIf
\EndFor
\State \Return $M$
\end{algorithmic}
\end{algorithm}

%% [#7 RESOLVIDA 12/11/2026] A nao reprodutibilidade da lista de candidatos fica
%% declarada como limitacao (os registros do SAGRES nao foram preservados).
O tamanho do conjunto candidato indicado pela equipe e os eventuais
descartes em cada critério não foram registrados durante a pesquisa; por isso,
não se pode afirmar quantos candidatos o filtro eliminou, e a etapa de
levantamento do conjunto candidato não é inteiramente reproduzível. Registra-se
essa limitação junto com outra, apontada também no artigo derivado: selecionar
módulos pela viabilidade de isolamento por \textit{stubs} sub-representa os
módulos com forte acoplamento ao \textit{hardware} ou ao escalonador, de modo
que a amostra não é estatisticamente representativa do PX4
(Seção~\ref{sec:ameacas-validade}).

A aplicação do procedimento ao PX4 Autopilot (ramo main, commit 8be2590c15) resultou em um conjunto de 12 módulos, organizados em três rodadas conforme descrito na Seção~\ref{sec:visao-geral}. Os quatro módulos da primeira rodada foram selecionados por sua criticidade direta para a segurança de voo: drivers de sensores (IMU BMI088 e GPS uBlox), \corr{comunicação (MAVLink, \textit{buffer} de transmissão)} e controle de entrada (Expo RC). As rodadas subsequentes ampliam a cobertura para monitoramento de sistema, navegação e estimação de estado.

%% LC 19 (resolvido 19/10/2026): citacao deSorbo2022 removida.
\corr{Foram priorizados componentes cuja falha pode resultar na perda de
controle do VANT ou na execução incorreta de missões autônomas.}
Os resultados quantitativos da verificação de cada módulo \revisao{(veredito, achados e tempo de execução)} são apresentados no Capítulo~\ref{cap:experimentos} (Experimentos e Resultados), juntamente com a discussão dos contraexemplos e das decisões de configuração do ESBMC específicas para cada módulo.

O driver IMU BMI088 integra acelerômetro e
giroscópio via SPI~\cite{PX4_Documentation}, realizando a leitura de
registradores, a conversão de dados brutos para unidades físicas (m/s² e rad/s),
 a aplicação de calibração e a leitura de temperatura. Dados incorretos da IMU
propagam-se para a estimação de atitude e podem causar instabilidade de
voo~\cite{Siddiqi2022}. O driver GPS uBlox realiza o parsing de mensagens do
protocolo UBX via UART~\cite{PX4_Documentation}; falhas no parser podem
corromper as coordenadas e comprometer funções de failsafe, como Return-to-Launch.
%% LC 20 (resolvido 19/10/2026): renomeado "MAVLink (buffer de transmissao)".
\corr{O módulo MAVLink implementa o protocolo de comunicação entre o autopiloto
e a estação de controle~\cite{Ficco2022, Hamza2024}. O \textit{harness} deste
trabalho não verifica o \textit{parser} de mensagens recebidas, e sim o
\textit{buffer} de transmissão (\texttt{send\_bytes}, em
\texttt{src/modules/mavlink/mavlink\_main.cpp}): o preenchimento não deve exceder
o tamanho máximo de pacote. Por isso, o módulo é referido neste texto como
``MAVLink (\textit{buffer} de transmissão)''.}
%% LC 21 (resolvido 19/10/2026): citacao trocada pelo arquivo-fonte.
\corr{O módulo Expo implementa a curva exponencial aplicada aos comandos do
controle remoto, com aritmética de ponto flutuante; a função verificada é
\texttt{math::expo()}, definida em
\texttt{src/lib/mathlib/math/Functions.hpp} no commit \texttt{8be2590c15}.}
Os módulos das rodadas seguintes ampliam a
cobertura para monitoramento, navegação e estimação de estado.

\begin{figure}[htb]
    \centering
    \resizebox{\textwidth}{!}{%
    \begin{tikzpicture}[
      font=\linespread{1}\footnotesize, >={Stealth[length=2mm]},
      ver/.style={draw=palEdgeDark, fill=palStrong, very thick, rounded corners=2pt,
        align=center, text width=3.2cm, minimum height=9mm},
      nver/.style={draw=palEdge, fill=palPaper, rounded corners=2pt, align=center,
        text width=3.2cm, minimum height=9mm},
      bus/.style={draw=palEdgeDark, very thick, fill=palFill, rounded corners=3pt, align=center,
        text width=2.6cm, minimum height=3.6cm},
      tit/.style={font=\small\bfseries}]
      \node[bus] (u) {\textbf{uORB}\\(barramento de mensagens\\\textit{publish--subscribe})};
      % publicadores
      \node[ver, left=22mm of u, yshift=24mm] (imu) {\textbf{IMU BMI088}};
      \node[ver, below=3mm of imu] (gps) {\textbf{GPS uBlox}};
      \node[ver, below=3mm of gps] (mag) {\textbf{Magnetômetro}\\e \textbf{Barômetro}};
      \node[ver, below=3mm of mag] (rc) {\textbf{RC Update}\\e \textbf{Expo RC}};
      \node[ver, below=3mm of rc] (bat) {\textbf{Battery Monitor}};
      \node[tit, above=3mm of imu] {Publicadores};
      % assinantes
      \node[ver, right=22mm of u, yshift=24mm] (ekf) {\textbf{EKF2}\\(covariância)};
      \node[ver, below=3mm of ekf] (nav) {\textbf{Navigator/Mission}};
      \node[nver, below=3mm of nav] (ctl) {Controle de posição\\e de atitude};
      \node[ver, below=3mm of ctl] (vtol) {\textbf{VTOL}};
      \node[ver, below=3mm of vtol] (pwm) {\textbf{PWM Output}};
      \node[tit, above=3mm of ekf] {Assinantes};
      \foreach \n in {imu,gps,mag,rc,bat} \draw[->] (\n.east) -- (u.west);
      \foreach \n in {ekf,nav,ctl,vtol,pwm} \draw[->] (u.east) -- (\n.west);
      \node[ver, below=10mm of u, text width=3.6cm] (mav) {\textbf{MAVLink}\\(enlace com a estação de solo)};
      \draw[<->] (u.south) -- (mav.north);
      \node[left=1mm of u.north west, anchor=south east, font=\scriptsize\itshape] {\texttt{orb\_publish()}};
      \node[right=1mm of u.north east, anchor=south west, font=\scriptsize\itshape] {\texttt{orb\_copy()}};
      % legenda
      \node[ver, text width=4mm, minimum height=4mm, below=6mm of mav.south west, xshift=-52mm, anchor=north west] (l1) {};
      \node[right=1mm of l1, anchor=west, font=\scriptsize] (t1) {módulo verificado neste trabalho (12 módulos)};
      \node[nver, text width=4mm, minimum height=4mm, right=6mm of t1] (l2) {};
      \node[right=1mm of l2, anchor=west, font=\scriptsize] {módulo não verificado};
    \end{tikzpicture}}
    \caption{Organização modular do PX4 Autopilot e interação entre módulos
    por meio do middleware uORB, com destaque (no tom mais escuro da paleta) para os doze módulos verificados
    \corr{na primeira campanha}. Elaborado pela autora com base
    em~\cite{PX4_Documentation}.}
    \label{fig:px4-modular}
\end{figure}

% =============================================================================
\section{Instrumentação e Preparação do Código}
\label{sec:instrumentacao}
% =============================================================================

A verificação formal de firmware embarcado apresenta desafios relacionados a
dependências de hardware. O PX4 utiliza intensivamente a HAL do NuttX e as bibliotecas de comunicação I2C, SPI e UART~\cite{PX4_Documentation}. Executar o verificador sobre o código original resultaria em falhas de compilação pela ausência dessas dependências no ambiente de verificação.

%% LC 22 (resolvido 19/10/2026): explicacao de nao determinismo removida daqui.
\corr{A Etapa~(ii) descreveu como essas dependências são abstraídas por valores
não determinísticos. Esta seção mostra como isso aparece nos \textit{harnesses}
da primeira campanha; a instrumentação da segunda campanha, com \textit{stubs}
de fato, é descrita na Seção~\ref{sec:segunda-campanha}.}

Nos \textit{harnesses} depositados, a abstração da E/S de
\textit{hardware} fica na própria função de teste: os bytes que o driver leria do
barramento (SPI, I2C ou UART) são declarados como valores não determinísticos e
passados diretamente ao trecho do driver sob verificação. A
Listagem~\ref{lst:stub_i2c} mostra a função \texttt{updateTemperature()} como
transcrita no \textit{harness} do IMU, e a Listagem~\ref{lst:main_imu}, a função
de teste que a exercita.

\begin{lstlisting}[
  caption={Função \texttt{updateTemperature()} transcrita no \textit{harness} do IMU BMI088 (\texttt{test\_bmi088\_imu.cpp}, comentários omitidos)},
  label={lst:stub_i2c},
  numbers=left,
  frame=single,
  language=C,
  breaklines=true,
  basicstyle=\ttfamily\small]
float updateTemperature(uint8_t temp_msb, uint8_t temp_lsb) {
    uint16_t Temp_uint11 = (temp_msb * 8) + (temp_lsb / 32);
    int16_t Temp_int11 = 0;

    if (Temp_uint11 > 1023) {
        Temp_int11 = Temp_uint11 - 2048;
    } else {
        Temp_int11 = Temp_uint11;
    }

    float temperature = (Temp_int11 * 0.125f) + 23.0f;
    return temperature;
}
\end{lstlisting}

\begin{lstlisting}[
  caption={Função de teste da temperatura no \textit{harness} do IMU BMI088 (\texttt{test\_bmi088\_imu.cpp}, comentários omitidos)},
  label={lst:main_imu},
  numbers=left,
  frame=single,
  language=C,
  breaklines=true,
  basicstyle=\ttfamily\small]
void test_temperature_calculation() {
    uint8_t temp_msb = nondet_uint8();
    uint8_t temp_lsb = nondet_uint8();

    float temperature = updateTemperature(temp_msb, temp_lsb);

    assert(temperature >= -40.0f && temperature <= 85.0f);

    assert(!isnan(temperature));
    assert(!isinf(temperature));
}
\end{lstlisting}

A função da Listagem~\ref{lst:stub_i2c} é uma transcrição do trecho
correspondente do driver BMI088 do PX4, copiada para dentro do
\textit{harness}, e não a função compilada a partir do código de produção. É assim
em todos os \textit{harnesses} \corr{da primeira campanha}: cada um extrai trechos do código-fonte do PX4,
que não foi modificado. Isso tem uma consequência que se registra aqui e nas
ameaças à validade (Seção~\ref{sec:ameacas-validade}): o que se verifica é a
cópia da lógica, e uma guarda presente no driver e omitida no trecho transcrito
produz um falso positivo --- exatamente o mecanismo do candidato do contador de
FIFO (Seção~\ref{subsec:imu-fifo}). É por isso que a conferência contra o código
de produção, na etapa~(v), é obrigatória. \corr{Essa limitação é a principal
motivação da segunda campanha.}

% =============================================================================
\section{Propriedades Verificadas}
\label{sec:propriedades}
% =============================================================================

As propriedades de segurança foram organizadas em duas categorias: automáticas,
verificadas pelo ESBMC a partir das flags de configuração, e customizadas,
expressas por assertions inseridas no código.

\subsection{Propriedades Automáticas}

O ESBMC verifica automaticamente diversas classes de defeitos quando habilitado
pelas flags de configuração. Seguindo a formalização de Cordeiro et al.~\cite{Cordeiro2012}, uma propriedade de safety $\phi$ é violada quando existe um estado alcançável $s_i$ tal que $s_i \not\models \phi$. \revisao{Em segurança de memória, o verificador detecta acessos fora dos limites de arranjos e desreferências de ponteiros nulos. Em aritmética inteira, identifica \textit{overflows} e \textit{underflows} de inteiros com sinal. Verifica ainda divisões por zero}~\cite{Menezes2024}.

A Tabela~\ref{tab:propriedades-automaticas} detalha as principais classes de defeitos verificadas automaticamente pelo ESBMC nesta pesquisa, suas descrições e as \textit{flags} de linha de comando associadas, conforme a documentação da ferramenta~\cite{Menezes2024}.

\begin{table}[!htb]
\centering
\caption{Propriedades automáticas verificadas pelo ESBMC neste trabalho}
\label{tab:propriedades-automaticas}
\footnotesize
\begin{tabular}{|p{3cm}|p{7cm}|p{3cm}|}
\hline
\textbf{Classe} & \textbf{Descrição} & \textbf{Flag ESBMC} \\
\hline
Acesso fora dos limites & Leitura/escrita em endereços inválidos de arrays e buffers & padrão (\texttt{--no-bounds-check} desliga) \\
\hline
Ponteiro nulo & Desreferência de ponteiro igual a \texttt{NULL} & padrão (\texttt{--no-pointer-check} desliga) \\
\hline
Overflow aritmético & Operações com inteiros com sinal que excedem a capacidade do tipo & \texttt{--overflow-check} \\
\hline
Divisão por zero & Operações de divisão ou módulo com divisor nulo & padrão (\texttt{--no-div-by-zero-check} desliga) \\
\hline
\corr{Deslocamento indefinido} & \corr{Deslocamento de bits com quantidade negativa ou maior que a largura do tipo (só na segunda campanha)} & \corr{\texttt{--ub-shift-check}} \\
\hline
\corr{Valores NaN} & \corr{Operação de ponto flutuante que produz NaN (só nos testes indicados da segunda campanha)} & \corr{\texttt{--nan-check}} \\
\hline
\end{tabular}

\footnotesize
\begin{flushleft}
\revisao{Acesso fora dos limites, ponteiro nulo e divisão por zero são verificados por padrão pelo ESBMC~8.1, que só oferece opções para desligá-los; a
verificação de \textit{overflow} aritmético é habilitada explicitamente na
invocação (Listagem~\ref{lst:comando-esbmc}).} \corr{As duas últimas linhas
valem apenas para a segunda campanha, com o ESBMC~8.5.}
\end{flushleft}
\end{table}

As verificações de NaN (\texttt{--nan-check}) e de vazamento de memória
(\texttt{--memory-leak-check}) não foram habilitadas na invocação da
Listagem~\ref{lst:comando-esbmc}, como no artigo derivado. O enunciado de que
os \textit{harnesses} de dez módulos não apresentam violações vale, portanto, para
as quatro classes da tabela e para as \textit{assertions} customizadas. NaN e
infinito são verificados apenas onde o \textit{harness} tem \textit{assertion}
explícita (\texttt{!isnan}, \texttt{!isinf}), o que ocorre em oito dos doze
módulos: IMU BMI088, Expo RC, EKF2, magnetômetro, barômetro, bateria, RC Update e
VTOL.

\revisao{Essas verificações são aplicadas uniformemente a todos os 12 módulos
analisados, sem necessidade de instrumentação específica por módulo. A flag
\texttt{--overflow-check} foi habilitada em todas as execuções, dado o uso
intensivo de aritmética inteira nos drivers de sensores do PX4.}
Ela não cobre o \textit{wraparound} de inteiros sem sinal, que a norma da
linguagem define e que só seria sinalizado com \texttt{--unsigned-overflow-check},
não usada nesta campanha, nem a conversão de um valor negativo para um tipo sem
sinal. Por isso, o \textit{underflow} do GPS não aparece como
\textit{overflow} aritmético, e sim como acesso fora dos limites dentro do
\texttt{memcpy} (Seção~\ref{sec:exemplo-ilustrativo}).

\subsection{Assertions Customizadas}

Além das verificações automáticas, foram adicionadas assertions baseadas em
especificações de fabricantes e requisitos do PX4~\cite{PX4_Documentation}.
Formalmente, cada assertion customizada expressa uma propriedade
$\phi_i$ da forma $\phi_i(s) \equiv \texttt{expr}$, onde \texttt{expr} é uma
condição booleana sobre o estado $s$ do programa~\cite{Cordeiro2012}.

Cada propriedade customizada é apresentada em três níveis de abstração: (i) \textbf{especificação em linguagem natural}, extraída do datasheet do fabricante ou da documentação do protocolo; (ii) \textbf{formalização}, expressando a propriedade como predicado lógico sobre o estado do programa; e (iii) \textbf{implementação}, por meio de \textit{assertion} inserida no \textit{harness} de verificação.

\textbf{Propriedade 1 --- Faixa operacional de temperatura do BMI088.}

\textit{Especificação em linguagem natural:} conforme o datasheet do fabricante Bosch Sensortec~\cite{Bosch2021}, a temperatura lida do sensor de temperatura integrado ao BMI088 deve estar contida na faixa operacional do componente, definida entre $-40\,^{\circ}$C e $+85\,^{\circ}$C. Valores fora dessa faixa indicam leitura inválida e não devem ser utilizados como entrada para os cálculos de compensação térmica do EKF2.

\textit{Formalização:}

\begin{equation}
\phi_{\text{temp}}(s) \equiv -40{,}0 \leq \texttt{temperature}(s) \leq 85{,}0
\label{eq:prop_temp}
\end{equation}

\revisao{Registre-se desde já que a passagem de uma faixa operacional de componente
físico para uma restrição sobre variável de programa é um passo de modelagem
contestável: o \textit{datasheet} declara em que faixa o sensor opera dentro da
especificação, e não que o software esteja proibido de representar uma leitura
fora dela. Foi precisamente esse passo que os mantenedores do PX4 recusaram, e a
discussão consta da Seção~\ref{subsec:imu-temperatura}. A propriedade é mantida
aqui como exemplo justamente por isso: ela ilustra a categoria de propriedade
cuja tradução exige mais cuidado, discutida na
Seção~\ref{subsec:proveniencia}.} Como a conversão afim do driver não
limita o valor, qualquer leitura bruta fora da faixa viola $\phi_{\text{temp}}$
por construção; por isso, como no artigo derivado, $\phi_{\text{temp}}$ é tratada
não como garantia do \textit{hardware}, mas como uma sonda de cobertura que
sinaliza a ausência de validação da faixa de entrada no driver.

\textit{Implementação:}

\begin{lstlisting}[
  numbers=left,
  frame=single,
  language=C,
  basicstyle=\ttfamily\small]
float temperature = updateTemperature(temp_msb, temp_lsb);
assert(temperature >= -40.0f && temperature <= 85.0f);
\end{lstlisting}

\textbf{Propriedade 2 --- Integridade do buffer de diagnóstico do GPS uBlox.}

\textit{Especificação em linguagem natural:} o nível de preenchimento do buffer de diagnóstico usado pela função \texttt{dumpGpsData()} deve permanecer dentro da capacidade do registro. No driver de produção, essa capacidade é \texttt{sizeof(dump\_data->data)}, 79 bytes (\texttt{msg/GpsDump.msg}); no \textit{harness}, ela é modelada pela macro \texttt{GPS\_DUMP\_DATA\_SIZE}, de valor 200, que não existe no código do PX4. Após qualquer operação de escrita no buffer, o comprimento acumulado \texttt{dump\_data.len} não deve exceder esse limite, sob pena de acesso fora dos limites da memória alocada.

\textit{Formalização:}

\begin{equation}
\phi_{\text{gps}}(s) \equiv \texttt{dump\_data.len}(s) \leq
\texttt{GPS\_DUMP\_DATA\_SIZE}
\label{eq:prop_gps}
\end{equation}

Como \texttt{len} é entrada livre, verificar $\phi_{\text{gps}}$ na entrada da
função seria trivialmente violado pela própria entrada. A verificação que importa
é a do acesso no ponto de uso, o \texttt{memcpy}: a violação reportada pelo ESBMC
é um acesso fora dos limites dentro do \texttt{memcpy} (\textit{dereference
failure}), e não a \textit{assertion} abaixo, que é a propriedade de mais alto
nível do \textit{harness}. \corr{Como discutido na
Seção~\ref{sec:exemplo-ilustrativo}, essa \textit{assertion} é tautológica na
forma depositada.}

\textit{Implementação:}

\begin{lstlisting}[
  numbers=left,
  frame=single,
  language=C,
  basicstyle=\ttfamily\small]
dumpGpsData(input_data, input_len, gps_dump_comm_mode_t::Full,
            /* msg_to_gps_device = */ true,
            &dump_buffer, gps_dump_comm_mode_t::Full);
assert(dump_buffer.len == 0 ||
       dump_buffer.len <= GPS_DUMP_DATA_SIZE ||
       (dump_buffer.len & 0x7F) <= GPS_DUMP_DATA_SIZE);
\end{lstlisting}

\textbf{\revisao{Propriedade 3 --- Ganho de Kalman no intervalo unitário.}}

\revisao{\textit{Especificação em linguagem natural:} no filtro de Kalman
implementado pelo módulo de fusão do magnetômetro, o ganho $K$ pondera quanto da
inovação da medição é incorporado à estimativa.} No \textit{harness}, ele
é calculado como $K = P/(P + R)$, em que $P$ é a variância do estado e $R$, a da
observação.  \revisao{Por definição, trata-se de um
fator de ponderação entre a predição e a observação, e portanto deve permanecer
no intervalo $[0;\,1]$. Valor fora desse intervalo indica erro numérico na
atualização da covariância e faz a estimativa divergir da observação em vez de
convergir para ela.}

\revisao{Diferentemente das duas anteriores, esta propriedade não deriva de
\textit{datasheet} nem de constante do código-fonte, mas da definição matemática
do estimador. É, por isso, a categoria de propriedade mais robusta das três:
independe de fidelidade entre documento e implementação e não envolve traduzir
grandeza física em variável de programa.}

\revisao{\textit{Formalização:}}

\begin{equation}
\phi_{\text{kalman}}(s) \equiv 0 \leq K(s) \leq 1
\label{eq:prop_kalman}
\end{equation}

\revisao{\textit{Implementação:}}

\begin{lstlisting}[
  numbers=left,
  frame=single,
  language=C,
  basicstyle=\ttfamily\small]
float compute_kalman_gain(float state_var, float obs_var)
{
    if (state_var + obs_var <= 0.0f) return 0.0f;
    return state_var / (state_var + obs_var);
}

float gain = compute_kalman_gain(state_var, obs_var);
assert(gain >= 0.0f);
assert(gain <= 1.0f);
assert(!isnan(gain));
\end{lstlisting}

As variâncias de entrada são não determinísticas, restritas por
\textit{assumptions} a valores finitos, não negativos e com soma positiva.

%% LC 24 (resolvido 19/10/2026): funcao escrita no harness, modela
%% VehicleMagnetometer.cpp.
\corr{A função \texttt{compute\_kalman\_gain} não é código do PX4: ela foi escrita
no próprio \textit{harness} (\texttt{test\_vehicle\_magnetometer.cpp}) para
modelar o cálculo do ganho feito em
\texttt{src/modules/sensors/vehicle\_magnetometer/VehicleMagnetometer.cpp}.
O veredito desta propriedade vale, portanto, para o modelo escrito no
\textit{harness}, e não diretamente para o código do PX4.}

No MAVLink, verificou-se que o preenchimento do buffer de transmissão
nunca excede sua capacidade, de 280 bytes, o tamanho máximo de pacote do
protocolo, mesmo com pacotes de entrada maiores que ela e após envios
sucessivos.

% =============================================================================
\subsection{\revisao{Proveniência das Propriedades}}
\label{subsec:proveniencia}

\revisao{A rastreabilidade entre cada propriedade verificada e o documento que a
originou é parte da metodologia: sem ela, não há como um terceiro auditar se a propriedade
traduz corretamente o requisito, nem como a comunidade avaliar um contraexemplo. A
Tabela~\ref{tab:proveniencia} registra essa correspondência para todas as propriedades
customizadas deste trabalho.}

%% LC 25 (resolvido parcialmente 19/10/2026): arquivo-fonte e funcao do harness
%% acrescentados; numeros de linha do PX4 e pagina do datasheet seguem pendentes.
%% LC 26 (resolvido 19/10/2026): Battery Monitor e Barometer acrescentados.
\begin{table}[!htb]
\centering
\caption{\revisao{Proveniência das propriedades customizadas verificadas} \corr{(primeira campanha)}}
\label{tab:proveniencia}
\footnotesize
\begin{tabular}{|p{3.0cm}|p{2.2cm}|p{2.5cm}|p{6.0cm}|}
\hline
\textbf{Propriedade} & \textbf{Módulo} & \textbf{Origem} & \textbf{Referência exata} \\
\hline
Faixa operacional de temperatura & IMU BMI088 & Datasheet do fabricante &
Bosch Sensortec, BST-BMI088-DS000-19, rev.~1.9~\cite{Bosch2021}; \corr{\textit{harness}: \texttt{test\_temperature\_calculation}} \\
\hline
Capacidade do buffer FIFO & IMU BMI088 & Código-fonte do PX4 &
Constante \texttt{FIFO::SIZE = 1024} em \texttt{Bosch\_BMI088\_ Accelerometer\_Registers.hpp}, \corr{l.~152}; \corr{\textit{harness}: \texttt{test\_fifo\_count\_calculation}} \\
\hline
Limite do buffer de diagnóstico & GPS uBlox & Código-fonte do PX4 &
\texttt{sizeof(dump\_data->data)}, 79 bytes, em \texttt{msg/GpsDump.msg}; no \textit{harness}, macro \texttt{GPS\_DUMP\_ DATA\_SIZE = 200}; \corr{\texttt{src/drivers/gps/gps.cpp}, l.~741} \\
\hline
Capacidade do buffer de transmissão & MAVLink & Especificação de protocolo &
Tamanho máximo de pacote MAVLink, 280 bytes (\texttt{MAVLINK\_MAX\_ PACKET\_LEN})~\cite{PX4_Documentation}; \corr{\texttt{Mavlink::send\_bytes} em \texttt{src/modules/mavlink/mavlink\_main.cpp}, l.~852} \\
\hline
Faixa de entrada do rádio & Expo RC, RC Update & Documentação do PX4 &
Intervalo normalizado $[-1;\,1]$~\cite{PX4_Documentation}; \corr{\texttt{math::expo()} em \texttt{src/lib/mathlib/math/Functions.hpp}, l.~84; \texttt{src/modules/rc\_update/rc\_update.cpp}} \\
\hline
Limites de PWM & PWM Output & Documentação do PX4 &
Faixas de motor e servo~\cite{PX4_Documentation}; \corr{\texttt{src/drivers/pwm\_out/PWMOut.cpp}; constantes \texttt{PWM\_MOTOR\_MIN}/\texttt{PWM\_MOTOR\_MAX} do \textit{harness}} \\
\hline
Variância não negativa & EKF2, Magnetômetro & Requisito matemático &
Propriedade de matriz de covariância; \corr{\texttt{src/modules/ekf2/EKF/covariance.cpp}} \\
\hline
Ganho de Kalman em $[0;\,1]$ & Magnetômetro & Requisito matemático &
Definição do filtro de Kalman; \corr{modelado no \textit{harness} a partir de \texttt{VehicleMagnetometer.cpp}} \\
\hline
Índice de missão válido & Navigator & Código-fonte do PX4 &
Limites de \textit{waypoint}; \corr{\texttt{src/modules/navigator/mission.cpp}} \\
\hline
\textit{Thrust} em $[0;\,0{,}9]$ & VTOL & Documentação do PX4 &
Limite de transição~\cite{PX4_Documentation}; \corr{\texttt{src/modules/vtol\_att\_control/vtol\_type.cpp}} \\
\hline
\corr{Fator de escala de tensão em $[1{,}0;\,1{,}3]$; avisos e falhas coerentes com os limiares} & \corr{Battery Monitor} & \corr{Código-fonte do PX4} & \corr{\texttt{src/lib/battery/battery.cpp} (\texttt{determineWarning}, \texttt{determineFaults}, \texttt{computeScale})} \\
\hline
\corr{Prioridade em $[0;\,100]$; pressão corrigida finita e em $[30000;\,120000]$ Pa} & \corr{Barômetro} & \corr{Código-fonte do PX4} & \corr{\texttt{src/lib/sensor\_calibration/Barometer.cpp}} \\
\hline
\end{tabular}
\end{table}

As \textit{assertions} de cada linha da tabela estão no
\textit{harness} do módulo correspondente, no repositório de artefatos, em
\texttt{resultados/<módulo>/test\_*.cpp}: \texttt{IMU\_BMI088},
\texttt{GPS}, \texttt{MAVLink}, \texttt{Expo\_RC}, \texttt{RC\_Update},
\texttt{PWM\_Output}, \texttt{EKF2\_Covariance}, \texttt{Magnetometer},
\texttt{Navigator\_Mission}, \texttt{VTOL}\corr{, \texttt{Battery\_Monitor} e
\texttt{Barometer}.}

\corr{No \textit{datasheet} do BMI088 (BST-BMI088-DS000-19, rev.~1.9, janeiro de 2024), a faixa operacional de $-40$ a $+85\,^{\circ}$C está na Seção~1.1, Tabela~1 (parâmetro \textit{Operating Temperature}, $T_A$), p.~7; a conversão do valor lido para graus Celsius, com resolução de $0{,}125\,^{\circ}$C por unidade, está na Seção~5.3.7 (registradores \texttt{0x22}--\texttt{0x23}), p.~28.}

\revisao{A distinção entre as três origens não é apenas documental: ela determina a
força da conclusão. Propriedades derivadas do \textbf{código-fonte} são as mais seguras, pois
a constante verificada é a mesma que o compilador usa. Propriedades derivadas da
\textbf{documentação ou de especificação de protocolo} dependem da fidelidade entre documento
e implementação. Propriedades derivadas de \textbf{datasheet} são as mais delicadas, porque o
documento descreve o comportamento do componente físico, e não necessariamente uma restrição
sobre a variável de software que o representa --- ponto discutido em detalhe na
Seção~\ref{subsec:imu-temperatura} do Capítulo~\ref{cap:experimentos}.}

\section{Análise de Contraexemplos}
\label{sec:contraexemplos}
% =============================================================================

%% LC 27 (resolvido 19/10/2026): secao reduzida a escala de classificacao.
%% LC 28 (resolvido 19/10/2026): falha reportada e limite do memcpy.
\corr{O procedimento de análise foi descrito na etapa~(v)
(Seção~\ref{subsec:etapa-replay}). Esta seção registra apenas a escala usada
para classificar os resultados, fixada no Capítulo~\ref{cap:introducao}:}
\textbf{violação alcançável no modelo},
quando a violação decorre dos estados admitidos pelas pré-condições do
\textit{harness} (caso do GPS); e
\textbf{falso positivo}, quando o candidato decorre de guarda presente no código
de produção e omitida no \textit{harness} (caso do contador de FIFO). A
observação de temperatura ficou fora dessas duas categorias: a violação é real
no modelo, mas a própria propriedade, derivada do \textit{datasheet}, foi
contestada pelos mantenedores. \corr{No caso do GPS, a falha reportada pelo
ESBMC não é a violação da \textit{assertion} da propriedade~\eqref{eq:prop_gps},
e sim a verificação automática de limites dentro do \texttt{memcpy}
(\textit{dereference failure}), consequência do \textit{underflow} descrito na
Seção~\ref{sec:exemplo-ilustrativo}.}

\corr{Na segunda campanha, essa mesma escala foi aplicada, com uma distinção que
a primeira já exigia: um \textbf{falso positivo de modelagem} decorre de uma
guarda de produção omitida no \textit{harness} (o contador de FIFO na primeira
campanha; o \texttt{bst} na segunda), ao passo que um \textbf{falso positivo da
ferramenta} decorre do verificador, e não do código. Os desfechos da segunda
campanha e os valores concretos de cada contraexemplo são apresentados no
Capítulo~\ref{cap:experimentos} e no Apêndice~\ref{apend:contraexemplos}.}

% =============================================================================
\section{\corr{Segunda Campanha: Código de Produção, \textit{Replay} e Sanitizadores}}
\label{sec:segunda-campanha}
% =============================================================================

\corr{A primeira campanha mostrou a limitação de verificar trechos transcritos:
o veredito vale para a cópia, e a conclusão sobre o driver depende de uma
conferência que, nos casos do FIFO e do GPS, inverteu o resultado. A segunda
campanha foi desenhada para reduzir essa distância. Ela mantém as oito etapas,
mas endurece a conferência contra o código de produção (etapa~v): cada candidato
alcançável é submetido a \textit{replay} do \emph{código real} do PX4 sob
sanitizadores dinâmicos (ASan/UBSan), e a validação cruzada com o CBMC passa a
cobrir todos os \textit{harnesses}, não apenas os candidatos. Dois achados
desta campanha foram reportados à comunidade em 12/11/2026, depois da revisão do
orientador: o do RoboClaw (\textit{Issue}~\#29785) e o do batmon
(\textit{Issue}~\#29787). Ambos aguardavam triagem na data de fechamento deste
texto.}

\corr{A campanha reúne dez módulos complementares aos doze da primeira, em três
frentes: (A)~oito candidatos adicionais em módulos de entrada e de sistema;
(B)~o driver de saída \textbf{RoboClaw}; e (C)~o driver de bateria
\textbf{batmon}. As frentes compartilham o protocolo de evidência --- verificação
com ESBMC e CBMC, \textit{replay} do código real e sanitizadores ---, mas
diferem no commit e no desfecho, como se detalha a seguir.}

\subsection{\corr{Frente A --- Oito candidatos adicionais}}
\label{subsec:2c-candidatos}

\corr{Oito módulos foram verificados no mesmo commit da primeira campanha
(\texttt{8be2590c15}), com ESBMC~8.1.0 e~8.5.0, CBMC~5.95.1 e os sanitizadores de
endereço (ASan) e de comportamento indefinido (UBSan). Um único script
(\texttt{rodar\_tudo.sh}) regenera todos os registros em qualquer máquina com as
três ferramentas, cada log com cabeçalho (data, máquina, pasta, comando e
versões) e código de saída diferente de zero se algum resultado divergir do
esperado. A Tabela~\ref{tab:candidatos-extras} resume os desfechos.}

\begin{table}[!htb]
\centering
\caption{\corr{Candidatos adicionais da segunda campanha (commit \texttt{8be2590c15})}}
\label{tab:candidatos-extras}
\scriptsize
\renewcommand{\arraystretch}{0.95}
\setlength{\tabcolsep}{3pt}
\begin{tabular}{|p{2.4cm}|p{5.2cm}|p{5.0cm}|}
\hline
\textbf{Módulo} & \textbf{Resultado no commit testado} & \textbf{Situação e desfecho} \\
\hline
\corr{\texttt{sumd.cpp}} & \corr{com quadro legítimo de 32 canais, \texttt{sumd\_decode()} escreve o 64.º byte em \texttt{sumd\_data[64]}, sobre \texttt{crc16\_high}} & \corr{\textbf{continua no \texttt{main}}; defeito reproduzido por \textit{replay} da função real; relatado por terceiro em 08/10/2026 (PR~\#29031, refeito como \#29117), fechado sem \textit{merge}; não reportado de novo} \\
\hline
\corr{\texttt{MicroStrain.cpp}} & \corr{construtor: \texttt{max\_len} pode ser 128 e \texttt{\_port[128]='\textbackslash0'} escreve fora do vetor de 128} & \corr{\textbf{continua no \texttt{main}}; defeito reproduzido (ASan); não reportado, pelo impacto desprezível (exige caminho de dispositivo com 128 caracteres ou mais)} \\
\hline
\corr{\texttt{externalChecks.cpp}} & \corr{\texttt{num\_events} (0--255) indexa \texttt{events[5]}} & \corr{era real; \textbf{corrigido \textit{upstream}} em \texttt{ca1d36476a} (PR~\#28598)} \\
\hline
\corr{\texttt{ft7\_technologies.cpp}} & \corr{contadores de byte indexam \texttt{\_readbuf[30]} sem limite} & \corr{era real; \textbf{corrigido \textit{upstream}} em \texttt{3acf7935a0} (PR~\#28600)} \\
\hline
\corr{\texttt{IridiumSBD.cpp}} & \corr{a guarda de espaço compara com \texttt{\_packet\_length}, mas o \texttt{memcpy} copia \texttt{buflen}} & \corr{era real; \textbf{corrigido \textit{upstream}} em \texttt{0c903d47d6} (PR~\#28599)} \\
\hline
\corr{\texttt{bst.cpp}} & \corr{violação apenas no \textit{harness} que omitia duas guardas já presentes no commit (de \texttt{e8e86a2e0f})} & \corr{\textbf{falso positivo de modelagem}; o \textit{harness} fiel verifica \textsc{successful}} \\
\hline
\corr{\texttt{logger.cpp}} & \corr{nenhuma violação de memória em \texttt{write\_info\_multiple()} para qualquer \texttt{key\_len}} & \corr{sem defeito de memória; o truncamento é alcançável, mas gera registro com \texttt{key\_len} errado, não escrita fora dos limites} \\
\hline
\corr{\texttt{Common.cpp}} & \corr{violação só com índice que a guarda de produção nunca permite} & \corr{\textbf{retirado}} \\
\hline
\end{tabular}
\end{table}

\corr{Dois pontos merecem registro. Primeiro, três candidatos
(\texttt{externalChecks}, \texttt{ft7} e \texttt{IridiumSBD}) eram violações reais
no commit testado e foram \emph{corrigidos pelo próprio PX4}, por outras pessoas,
em setembro de 2026; as correções \textit{upstream} corroboram, de forma
independente, que os achados eram defeitos reais, e por isso nenhuma \textit{issue}
foi aberta para eles. O \texttt{sumd} teve trajetória parecida: o mesmo defeito
foi relatado por outra pessoa, mas o PR foi fechado sem \textit{merge}, e o código
permanece igual no \texttt{main}. Segundo, a verificação com o CBMC usou as mesmas
pré-condições e limites de desdobramento do ESBMC, com \textit{unwinding
assertions} e \textit{bounds} que cobrem os laços (por exemplo, 260 no
\texttt{externalChecks} e 345 no \texttt{IridiumSBD}); as duas ferramentas
acusam as mesmas violações. O caso do \texttt{sumd} precisou de
\texttt{--unwind 70}; como o modo \texttt{--multi-property} do ESBMC ultrapassa
quinze minutos nesse limite, a contagem de propriedades foi tomada do CBMC.}

\subsection{\corr{Frente B --- Driver de saída RoboClaw}}
\label{subsec:2c-roboclaw}

\corr{O driver \texttt{roboclaw} converte os comandos de controle em pacotes do
protocolo \textit{Packet Serial} enviados por UART aos controladores de motor.
Dois defeitos foram verificados no commit \texttt{8be2590c15}, com ESBMC~8.2.0
(solver Z3) e CBMC~5.95.1 (Tabela~\ref{tab:roboclaw-resultados}).}

\begin{table}[!htb]
\centering
\caption{\corr{RoboClaw: verificação formal dos dois defeitos}}
\label{tab:roboclaw-resultados}
\footnotesize
\begin{tabular}{|p{7.6cm}|p{2.6cm}|p{2.6cm}|}
\hline
\textbf{Defeito} & \textbf{Sem correção} & \textbf{Com \texttt{-DCORRECAO}} \\
\hline
\corr{Bug 1 --- byte de velocidade $>127$ em \texttt{sendUnsigned7Bit()}} & \corr{\textsc{failed}} & \corr{\textsc{successful}} \\
\hline
\corr{Bug 2 --- erro de \texttt{write()} não detectado em \texttt{writeCommandWithPayload()}} & \corr{\textsc{failed}} & \corr{\textsc{successful}} \\
\hline
\end{tabular}

\scriptsize
\begin{flushleft}
\corr{Mesmo padrão nas duas ferramentas (ESBMC e CBMC), na máquina da autora, em
28/10/2026. Registros e \textit{hashes} em
\texttt{testes\_hardware/roboclaw/evidencia\_maquina/}.}
\end{flushleft}
\end{table}

\corr{O \textbf{Bug~1} decorre do caminho de terminação de voo: o
\texttt{module.yaml} admite valores de \textit{failsafe} até 257 e o
\textit{mixer} os repassa ao driver sem limitar, enquanto o manual da RoboClaw
(Rev.~6.0, pp.~48--49) restringe o byte de dados dos comandos 0, 1, 4 e 5 à faixa
$0$--$127$. Um valor de \textit{failsafe} igual a 257 produz o byte 130. O
\textit{replay} com UART emulada, sobre a lógica real do driver, envia de fato o
byte 130 de forma determinística. O \textbf{Bug~2} vem de armazenar o retorno de
\texttt{write()} num \texttt{size\_t}: quando \texttt{write()} devolve $-1$, o valor
vira um inteiro sem sinal grande, a checagem \texttt{bytes\_written < 2} não
detecta o erro e a função retorna \texttt{OK}; o driver então espera uma resposta
que não virá, e a falha só aparece como \textit{timeout} em \texttt{select()}.
A correção dos dois está em \texttt{correcao\_roboclaw.patch}.}

\corr{O driver foi reescrito no \texttt{main} do PX4 em 28/10/2026
(PR~\#29473), depois do commit analisado. Pela leitura do código novo, o limite de
\textit{failsafe} passou a 254, o valor é restringido antes da conversão e o
retorno da escrita é guardado em \texttt{ssize\_t}; essa versão não foi
verificada nem ensaiada neste trabalho. O código antigo continua no ramo
\texttt{release/1.18}, em que \texttt{Roboclaw.cpp}, \texttt{Roboclaw.hpp} e
\texttt{module.yaml} são idênticos aos do commit analisado (conferido em
12/11/2026, commit \texttt{8303efda43cb}). Por isso os dois defeitos foram
reportados com esse ramo como alvo (\textit{Issue}~\#29785).}

\corr{\paragraph{Validação em hardware do Bug~1.} O Bug~1 foi observado também
em bancada, numa Pixhawk~6X~RT com firmware compilado do mesmo commit
(\texttt{px4\_fmu-v6xrt\_default}). A porta TELEM2 foi ligada a um ESP32 que
faz o papel da RoboClaw: responde às consultas do driver e registra cada byte
recebido. Com \texttt{RBCLW\_FAIL1} igual a 257 e a terminação de voo acionada
pelo console (\texttt{commander termination}), a placa passou a transmitir o
pacote \texttt{80 00 82 8A 90}: endereço \texttt{0x80}, comando~0, byte de
dados 130 e CRC válido. O pacote se repetiu 24\,136 vezes nas duas capturas
guardadas, sem erro de CRC. Dois ensaios complementares, na mesma montagem,
isolam a causa (Tabela~\ref{tab:roboclaw-hardware}): com o parâmetro em 254, o
maior valor que ainda resulta em byte válido, o mesmo firmware envia 127; e,
com a correção aplicada e o parâmetro de volta em 257, o byte também fica em
127, em 5\,267 pacotes e nenhum acima da faixa.}

\begin{table}[htb]
\centering
\caption{\corr{RoboClaw: byte de dados observado na serial da Pixhawk~6X~RT durante a terminação}}
\label{tab:roboclaw-hardware}
\begin{tabular}{llll}
\hline
\corr{\textbf{Firmware}} & \corr{\textbf{\texttt{RBCLW\_FAIL1}}} & \corr{\textbf{Pacote}} & \corr{\textbf{Byte de dados}} \\
\hline
\corr{Original} & \corr{257} & \corr{\texttt{80 00 82 8A 90}} & \corr{130 (fora de 0--127)} \\
\corr{Original} & \corr{254} & \corr{\texttt{80 00 7F B4 22}} & \corr{127} \\
\corr{Corrigido} & \corr{257} & \corr{\texttt{80 00 7F B4 22}} & \corr{127} \\
\hline
\end{tabular}
\end{table}

\corr{Esse resultado tem quatro limites, que delimitam o que ele prova.
Primeiro, a RoboClaw é emulada: o ensaio mostra que a placa \emph{transmite} o
byte inválido, mas não o que o controlador real faz ao recebê-lo, o que o
manual não especifica. Segundo, o firmware não é o de fábrica: o driver não
integra a configuração padrão dessa placa e foi incluído com
\texttt{CONFIG\_DRIVERS\_ROBOCLAW=y}; além disso, foi preciso inicializar dois
membros da classe (\texttt{\_uart\_fd\_set} e \texttt{\_uart\_fd\_timeout}).
Sem isso, em duas partidas registradas, o driver não emitiu mensagem alguma e a
fila \texttt{wq:hp\_default} ficou quase parada (\texttt{battery\_status} a
$0{,}1$~Hz, contra $99{,}9$~Hz depois da alteração; \texttt{roboclaw} a
$0{,}0$~Hz, contra $3{,}4$~Hz). Pelo código, \texttt{Run()} chama a atualização
do \textit{mixer} antes de inicializar a UART, de modo que o primeiro envio
alcança \texttt{select()} com esses membros ainda não inicializados; a causa não
foi confirmada com depurador, e os parâmetros do driver estavam sendo ajustados
entre essas partidas. Registra-se, por isso, como observação de bancada, e não
como terceiro defeito confirmado. Essa alteração não toca o caminho de saída. Terceiro, a placa registrou \textit{timeout} de ACK
durante os ensaios, de modo que o driver foi observado no caminho de erro de
resposta; a medida, que é o byte lido diretamente do fio, não depende disso.
Quarto, só o Bug~1 foi exercitado em hardware; a evidência do Bug~2 continua
sendo a formal e o \textit{replay} com UART emulada. As capturas, os
\textit{diffs} dos firmwares e o código do emulador estão em
\texttt{testes\_hardware/roboclaw/hardware\_pixhawk6xrt/}.}

\subsection{\corr{Frente C --- Driver de bateria batmon}}
\label{subsec:2c-batmon}

\corr{No driver \texttt{smart\_battery/batmon}, o número de células lido do
registrador SMBus da bateria define \texttt{\_cell\_count} até
\texttt{MAX\_CELL\_COUNT}${}=16$, e o laço de publicação escreve
\texttt{new\_report.voltage\_cell\_v[i]} para \texttt{i < \_cell\_count}. Mas
\texttt{voltage\_cell\_v} é um \texttt{float32[14]}
(\texttt{msg/versioned/BatteryStatus.msg}). Uma bateria legítima de 15S ou 16S,
dentro dos registradores que o próprio driver define, já basta para escrever fora
do vetor. Com 16 células, o valor extra cai em \texttt{full\_charge\_capacity\_wh}
(o \textit{replay} publica 3,715, a tensão da 16.ª célula), e \texttt{cell\_count}${}=16$
é publicado num vetor de 14 posições. O driver está nas configurações padrão de
várias placas (fmu-v3 a v6u, cubeorange, entre outras).}

\corr{A Tabela~\ref{tab:batmon-resultados} reúne os resultados. Diferentemente das
outras frentes, a conferência foi feita contra o \texttt{main} de 13/10/2026
(commit \texttt{c9b9e4e25e}), e não \texttt{8be2590c15} --- ressalva que se
mantém explícita.}

\begin{table}[!htb]
\centering
\caption{\corr{batmon: resultados (commit \texttt{c9b9e4e25e})}}
\label{tab:batmon-resultados}
\footnotesize
\begin{tabular}{|p{6.8cm}|p{3.2cm}|p{2.6cm}|}
\hline
\textbf{Ferramenta / verificação} & \textbf{Código atual} & \textbf{\texttt{-DCORRECAO}} \\
\hline
\corr{ESBMC~8.1.0 (\texttt{--unwind 20 --overflow-check})} & \corr{\textsc{failed}} & \corr{\textsc{successful}} \\
\hline
\corr{ESBMC~8.5.0 (\texttt{--multi-property})} & \corr{\textsc{failed} (2 de 13: \texttt{i < 14} e \textit{array bounds})} & \corr{\textsc{successful}} \\
\hline
\corr{CBMC~5.95.1 (\texttt{--unwind 20}, \textit{bounds} e \textit{pointer check})} & \corr{1 de 15 falha} & \corr{0 de 15} \\
\hline
\corr{UBSan (\texttt{./replay 16})} & \corr{\texttt{index 14 out of bounds for type 'float [14]'}} & \corr{---} \\
\hline
\corr{\textit{replay} sem sanitizador} & \corr{\texttt{full\_charge\_capacity\_wh=3.7150}} & \corr{---} \\
\hline
\end{tabular}
\end{table}

\corr{O mesmo defeito foi corrigido no driver UAVCAN/DroneCAN
(\texttt{cbat\_sub\_cb}) pelo PR~\#28606, em 14/09/2026, relatado por um terceiro;
o \texttt{batmon} ficou de fora dessa correção, e as buscas por ``batmon'',
``cell\_count'' e ``voltage\_cell\_v'' nas \textit{issues} do PX4 não encontraram
registro. O código foi reconferido em 12/11/2026 contra o \texttt{main} de
11/11/2026 (commit \texttt{8e94784e1dc0}), sem alteração no driver nem na
mensagem, e o \textit{replay} foi reexecutado com o mesmo resultado. O achado foi
então reportado (\textit{Issue}~\#29787). Não houve ensaio em hardware para este
driver.}

\subsection{\corr{Protocolo de conferência contra o código de produção}}
\label{subsec:2c-protocolo}

\corr{As três frentes seguem o mesmo protocolo de evidência, que operacionaliza a
etapa~(v) (Seção~\ref{subsec:etapa-replay}). As entradas do código verificado são
declaradas não determinísticas e restritas por pré-condições
(\texttt{\_\_VERIFIER\_assume}, aceita pelo ESBMC e pelo CBMC), com a origem de
cada pré-condição anotada no cabeçalho do \textit{harness}: um limite de
parâmetro declarado no \texttt{module.yaml}, o modo como o próprio PX4 chama a
função, ou uma constante do código. No RoboClaw, por exemplo, a saída recebida do
\textit{mixer} foi restrita à faixa $[0;\,257]$ porque o \texttt{module.yaml}
declara \texttt{failsafe: \{min: 0, max: 257\}}; uma pré-condição mais estreita
teria ocultado o Bug~1.}

\corr{Cada candidato alcançável passa por \textit{replay}: a função \emph{real}
do PX4 é compilada com um compilador comum (g++) e executada com os valores do
contraexemplo, sob ASan e UBSan. O critério é binário --- se o código real
reproduz o efeito previsto (o byte 130 no RoboClaw, o acesso a
\texttt{voltage\_cell\_v[14]} no batmon, a escrita em \texttt{\_port[128]} no
MicroStrain), o candidato é defeito reproduzido; se não reproduz, investiga-se a
guarda de produção que o \textit{harness} omitiu, como no \texttt{bst}. Quando há
correção, o mesmo \textit{harness} é reverificado com o \textit{patch}
(\texttt{-DCORRECAO}) e deve passar; para os três candidatos corrigidos
\textit{upstream}, a reverificação usa a correção do próprio PX4
(\texttt{-DCORRECAO\_UPSTREAM}).}

\corr{A validação cruzada com o CBMC foi estendida a todos os \textit{harnesses},
com as opções mapeadas para os equivalentes mais próximos
(Tabela~\ref{tab:equivalencia-opcoes}). Divergências entre as ferramentas são
registradas como resultado da comparação: elas vêm de modelagem de biblioteca
(por exemplo, \texttt{isfinite} tratado pelo CBMC como retorno arbitrário, que
exigiu um prelúdio de modelagem) ou de diferença de versão na contagem de
propriedades, e não de desacordo sobre a existência das violações.}

\begin{table}[!htb]
\centering
\caption{\corr{Equivalência de opções entre ESBMC e CBMC~5.95}}
\label{tab:equivalencia-opcoes}
\footnotesize
\begin{tabular}{|p{6.5cm}|p{7.0cm}|}
\hline
\textbf{ESBMC} & \textbf{CBMC} \\
\hline
\corr{checagens padrão (limites, ponteiros, divisão por zero)} & \corr{\texttt{--bounds-check --pointer-check --div-by-zero-check}} \\
\hline
\corr{\texttt{--overflow-check}} & \corr{\texttt{--signed-overflow-check --float-overflow-check}} \\
\hline
\corr{com \textit{unwinding assertions} (padrão) / \texttt{--no-unwinding-assertions}} & \corr{\texttt{--unwinding-assertions} / padrão} \\
\hline
\corr{\textit{solver} Boolector ou Z3} & \corr{MiniSat (padrão)} \\
\hline
\end{tabular}
\end{table}

\corr{A segunda campanha rendeu ainda um relato fora do PX4. Ao verificar o
\texttt{logger}, o modelo de \texttt{snprintf} do ESBMC devolvia um comprimento
incorreto quando o argumento de \texttt{\%s} era montado em tempo de execução,
o que produzia alarme falso num programa correto. O caso foi reduzido a um
programa mínimo e reportado ao projeto ESBMC (\textit{Issue}~\#8524, de
12/10/2026); os mantenedores confirmaram o problema também no ramo principal e o
corrigiram (PR~\#8563).}

% =============================================================================
\section{Resumo}
\label{sec:resumo-metodologia}

Este capítulo apresentou a metodologia de verificação formal aplicada ao PX4 Autopilot com o ESBMC. O procedimento foi organizado em oito etapas e aplicado em duas campanhas, que abrangeram 22 módulos distintos. Foram descritos os critérios de seleção dos módulos, a abstração das dependências de \textit{hardware} e do NuttX, a definição das propriedades de segurança e a configuração do verificador. Também foram apresentados os procedimentos de análise dos contraexemplos, conferência com o código de produção, validação cruzada com o CBMC e verificação de correções. A aplicação ao GPS uBlox ilustrou o fluxo da primeira campanha; a segunda ampliou a análise para dez módulos complementares e incorporou a execução do código original sob sanitizadores. Os resultados quantitativos, bem como a classificação e a discussão dos achados, são apresentados no Capítulo~\ref{cap:experimentos}.
